% Purpose: Mathematical framework, stochastic point-process foundations,
% self-exciting Hawkes/ETAS triggering, point-adaptive density estimation,
% vectorized ADP++ clustering, and topological cluster validity (PAk-DS Score)
% for seismic sequences and complex non-convex manifolds.
% Status: Under Review; mathematical audit incorporated, scientific
% review pending.
% Source-of-truth: OGS/src/ogsclustering.py, OGS/test/testogsclustering.py,
% doc/UNITS/supplementary/Clustering.tex, doc/UNITS/lib.bib

\providecommand{\bm}[1]{\mathbf{#1}}
\providecommand{\coloneqq}{\mathrel{\mathop:}=}
\providecommand{\mathds}[1]{\mathbb{#1}}
\providecommand{\pak}{\textsc{PAk}}
\providecommand{\kstar}{k^{*}}
\providecommand{\xvec}{\bm{x}}
\providecommand{\mdl}{\mathrm{MDL}}
\providecommand{\Zscore}{Z}
\providecommand{\etas}{\textsc{etas}}
\providecommand{\Var}{\operatorname{Var}}
\providecommand{\Cov}{\operatorname{Cov}}

\label{chap:clustering}

The identification, characterisation, and spatiotemporal
delimitation of seismic
clusters represent foundational challenges in quantitative seismology, hazard
assessment, and statistical physics. While tectonic seismicity is driven over
geological epochs by secular tectonic loading, the observed event catalog is
dominated by intense spatiotemporal clustering: coseismic elastostatic stress
transfer generates power-law decaying aftershock cascades, whereas transient
hydrothermal pressurisation, magmatic intrusions, and slow-slip
transients drive
complex earthquake swarms.

Disentangling these superimposed physical processes requires moving beyond
heuristic window-based declustering and Euclidean distance metrics.
This chapter
develops a unified, mathematically consistent framework that synthesises
measure-theoretic stochastic point processes, self-exciting marked Hawkes and
\ac{ETAS} mechanics, non-parametric
Point-Adaptive $k$-Nearest Neighbour (\pak{}) density estimation, vectorized
Advanced Density Peaks (\textsc{ADP++}) clustering, and topological cluster
validity via the \pak{} Density Separation Score (DS-Score).

The inspected executable evidence concerns density estimation, ADP/ADP++
clustering, and PAk-DS scoring. The temporal unified model and its EM/MDL
workflow are proposals, not verified implementations. Exact distributional
identities refer to explicitly stated ideal sampling experiments; adaptive
selection and physical attribution do not inherit them automatically.

\paragraph{Epistemic Governance and Document Blueprint.}
To maintain strict evidentiary boundaries, every mathematical derivation,
computational algorithm, and empirical evaluation in this chapter is classified
according to the six-tier epistemic taxonomy established in
\Cref{tab:clustering_epistemic_matrix}.

\begin{table}[htbp]
  \centering
  \small
  \begin{threeparttable}
    \caption{Evidence classes used in \cref{chap:clustering}.
      Inputs, constitutive models, source operations, and proposed
      interpretations carry different verification requirements.
      This audit inspected clustering code and tests, not catalogue
    measurements or benchmark run records.}
    \label{tab:clustering_epistemic_matrix}
    \begin{tabularx}{\linewidth}{l >{\raggedright\arraybackslash}X
      >{\raggedright\arraybackslash}X >{\raggedright\arraybackslash}X}
      \toprule
      \textbf{Tier} & \textbf{Epistemic Scope} & \textbf{Artifact /
      Formulation} & \textbf{Verification Anchor} \\
      \midrule
      \textbf{Level 1} & Empirical inputs & Arrival times,
      locations, magnitudes & Not reanalysed here \\
      \textbf{Level 2} & Constitutive models & Stress change,
      pressure diffusion & Stated physical assumptions \\
      \textbf{Level 3} & Deterministic computation & Counts,
      volumes, pointer updates & Source and test inspection \\
      \textbf{Level 4} & Model estimates & Dimension, adaptive
      density, temporal rates & Estimation assumptions;
      unified EM is proposed \\
      \textbf{Level 5} & Theoretical baselines & Poisson laws,
      information, persistence & Explicit sampling and
      regularity conditions \\
      \textbf{Level 6} & Interpretations & Triggering causes,
      thermodynamic analogies & Require independent evidence \\
      \bottomrule
    \end{tabularx}
    \begin{tablenotes}[flushleft]
      \footnotesize
    \item [Analyst Classification] These levels organize evidence,
      not degrees of truth. EM denotes expectation--maximization.
      No test pass rate, empirical result, or physical attribution
      is established by this table.
    \end{tablenotes}
  \end{threeparttable}
\end{table}

This chapter is organised into nine logically integrated sections:
\begin{itemize}[noitemsep,topsep=3pt]
  \item \Cref{sec:seismic_clustering_problem} examines the
    physical phenomenology of seismic sequences and details the
    mathematical failure modes of classical declustering and
    heuristic clustering metrics.
  \item \Cref{sec:point_process_foundations} derives the
    measure-theoretic foundations of stochastic point processes,
    \ac{IPP}, their grand canonical ensemble isomorphism, and optimal
    piecewise-constant changepoint segmentation via \ac{PELT}.
  \item \Cref{sec:pak_theory} establishes the differential
    geometry and non-parametric statistics of the \pak{} framework,
    including \ac{TWO-NN} intrinsic dimension estimation, adaptive
    neighborhood selection $\kstar$ via sequential likelihood-ratio
    tests, and closed-form variance stabilization.
  \item \Cref{sec:clustering_hawkes_etas} derives marked
    self-exciting Hawkes and \ac{ETAS} processes, proves the
    analytical branching ratio stability theorem, and formulates the
    constitutive mechanics of Coulomb stress transfer and
    pore-pressure diffusion.
  \item \Cref{sec:info_thermo_unified} formalises the
    thermodynamic free-energy landscape, establishes the
    information-geometric Euclidean flatness in log-density
    coordinates, and formulates the multi-scale Unified Generative
    Model with joint \ac{MDL}--\pak{} model selection.
  \item \Cref{sec:adppp_optimizations} presents the
    vectorized \textsc{ADP++} clustering engine, providing fixed-forest
    convergence proofs and scoped complexity bounds for fused
    union-find, saddle-point detection, $k$-peaks, and multimodality merging.
  \item \Cref{sec:ds_score_persistence} formalises the \pak{}
    Density Separation Score, relates its contrasts to ideal 0-homology
    persistence, and identifies the limits of the Kramers analogy and
    comparisons with classical validity indices.
  \item \Cref{sec:clustering_benchmarks} evaluates the
    complete unified framework across canonical synthetic 2D tests,
    a 21-manifold non-convex benchmark suite, and realistic
    synthetic seismic sequence scenarios.
  \item \Cref{sec:clustering_discussion} discusses boundary
    conditions, spatial extensions to $(x,y,z,t)$, and operational
    streaming deployment.
\end{itemize}

% ==============================================================================
% SECTION 1: THE SEISMIC CLUSTERING PROBLEM
% ==============================================================================
\section{The Seismic Clustering Problem}
\label{sec:seismic_clustering_problem}

\subsection{Phenomenology of Earthquake Sequences}
\label{subsec:seismic_phenomenology}

In observational seismology, earthquake catalogues comprise discrete
space-time-magnitude tuples recorded across regional sensor networks
(\textbf{Level~1: Empirical Observations}). When examined across
temporal windows ranging from hours to decades, seismic activity
exhibits pronounced departures from stationarity. Physically,
tectonic seismicity organizes into three dominant phenomenological regimes:

\begin{enumerate}[label=(\alph*),leftmargin=2em]
  \item \textbf{Secular Background Seismicity:} Driven by deep,
    steady tectonic plate motion and ductile shear at depth. In the
    absence of localized stress perturbations, background events
    occur as uncorrelated spatial Poisson arrivals \parencite{gardner1974}
    with constant or slowly varying rate $\mu_0 > 0$.
  \item \textbf{Mainshock--Aftershock Cascades:} Triggered by the
    abrupt elastostatic and dynamic stress drop of a major rupture
    ($\Delta\bm{\sigma} \sim \qty{1}{\mega\pascal}$ to
    $\qty{10}{\mega\pascal}$). The sequence initiates with an
    exceptional event that dominates the seismic moment budget of
    the sequence ($M_0^{(\mathrm{main})} \ge \sum_j M_0^{(j)}$),
    obeying B\r{a}th's empirical law \parencite{bath1965,richter1958},
    wherein the magnitude differential between the mainshock and
    the largest triggered aftershock is approximately constant
    \parencite{shcherbakov2004,helmstetter2003bath,console2003}:
    \begin{equation}
      \Delta m \coloneqq m_{\mathrm{main}} -
      m_{\mathrm{after}}^{\max} \approx \num{1.2}.
      \label{eq:baths_law}
    \end{equation}
    The subsequent relaxation rate decays monotonically following
    the modified Omori--Utsu law $\lambda(t) \propto (t + c)^{-p}$
    \parencite{omori1894,utsu1961,utsu1995} with exponent
    $p \approx \numrange{1.0}{1.3}$, while magnitude frequencies conform
    to the baseline Gutenberg--Richter distribution
    \parencite{ishimoto1939,gutenberg1944} with $b \approx
    \numrange{0.8}{1.1}$.
  \item \textbf{Earthquake Swarms:} Characterized by a protracted
    sequence of events of comparable magnitudes without an
    identifiable initial mainshock ($\Delta m \le \num{0.5}$)
    \parencite{mogi1963,sykes1970,vidale2006}. Swarms
    exhibit emergent or accelerating initiation, irregular
    multi-peak burst rates, and strong spatiotemporal hypocenter
    migration. Swarms are physically driven by non-seismic forcing
    mechanisms: pore-fluid overpressure diffusion along permeable
    fault damage zones ($\partial p/\partial t = D\nabla^2 p \equiv
    D\operatorname{div}(\operatorname{grad} p)$)
    \parencite{nur1972,shapiro1997,shapiro2002}, magmatic intrusions,
    or aseismic slow fault slip transients \parencite{hainzl2005}.
    Their frequency-magnitude distribution frequently exhibits
    anomalously high $b$-values ($b \in [\num{1.3}, \num{2.5}]$), reflecting
    localized shear stress reduction under elevated pore-fluid pressures
    \parencite{scholz1968,wiemer2002}.
\end{enumerate}

\subsection{Pathologies of Classical Declustering Algorithms}
\label{subsec:declustering_limitations}

Historically, \ac{PSHA} require isolating independent background
events by removing clustered aftershocks—a procedure termed \emph{catalogue
declustering}. Classical windowing methods, most notably the
algorithms of \citet{gardner1974} and \citet{reasenberg1985}, impose
deterministic spatial and temporal interaction envelopes around
identified triggers:
\begin{equation}
  \Delta t \le T(M), \qquad \Delta r \le L(M),
  \label{eq:classical_window}
\end{equation}
where $T(M)$ and $L(M)$ are empirical step functions or power-law
scalings of parent magnitude.

These heuristic windowing approaches exhibit severe theoretical and
operational deficiencies:
\begin{itemize}[noitemsep,topsep=2pt]
  \item \textbf{Arbitrary Binary Classification:} A sharp
    spatiotemporal boundary forces a hard binary assignment (cluster
    member vs.\ independent background), ignoring the fundamentally
    continuous, probabilistic nature of stress-mediated triggering.
  \item \textbf{Chaining and Cluster Distortion:} In high-density
    seismotectonic domains, adjacent sequences overlap. Sequential
    windowing causes runaway chaining, aggregating distinct,
    physically unrelated swarms and background seismicity into
    single, distorted macro-clusters.
  \item \textbf{Erasure of Swarm Dynamics:} Because window sizes
    scale strictly with magnitude, swarm sequences containing
    hundreds of low-to-moderate magnitude events ($m \le \num{3.5}$) are
    either fragmented into dozens of spurious micro-sequences or
    entirely filtered out as background noise, destroying the
    continuous rate and migration signatures vital for tracking
    fluid diffusion.
  \item \textbf{Regional Bias:} Empirical window thresholds
    calibrated for strike-slip regimes (e.g., Southern California)
    fail catastrophically when applied to complex collisional thrust
    belts, subduction zones, or volcanic hydrothermal systems.
\end{itemize}
Modern stochastic and topological frameworks address these failure modes
by replacing rigid windows with continuous probabilistic triggering models
\parencite{zhuang2002,marsan2008} or rescaled nearest-neighbour distances
in space--time--magnitude manifolds
\parencite{baiesi2004,zaliapin2008,zaliapin2013}.

\subsection{Failure of Classical Clustering Metrics on Complex Manifolds}
\label{subsec:classical_metrics_failure}

When unsupervised machine learning algorithms are tasked with
partitioning earthquake catalogs or spatial hypocentral clouds,
evaluating the resulting partitions requires \acp{CVI}. The most widely adopted
validation metrics \eg{the \ac{SC}, the \ac{CH}, and the \ac{DB}}
are constructed under the implicit geometric assumption that
clusters form convex, isotropic, hyperspherical distributions in
Euclidean space:

\begin{equation}
  S_{\mathrm{Sil}} \coloneqq \frac{1}{N}\sum_{i=1}^N \frac{b(i) -
  a(i)}{\max(a(i), b(i))}, \qquad
  \mathrm{CH} \coloneqq \frac{\mathrm{Tr}(\mathbf{B})/(K -
  1)}{\mathrm{Tr}(\mathbf{W})/(N - K)}, \qquad
  \mathrm{DB} \coloneqq \frac{1}{K}\sum_{k=1}^K \max_{l \ne k} \left(
  \frac{s_k + s_l}{d(\mathbf{c}_k, \mathbf{c}_l)} \right),
  \label{eq:classical_cvis}
\end{equation}
where $\mathbf{B}$ and $\mathbf{W}$ are the between-cluster and
within-cluster scatter matrices, $\mathbf{c}_k$ are cluster
centroids, and $s_k=|C_k|^{-1}\sum_{i\in C_k}\|\mathbf{x}_i-\mathbf{c}_k\|$
are mean centroid distances, matching scikit-learn. Here $a(i)$ is the
mean distance to other members of the assigned cluster and $b(i)$ the
smallest mean distance to another cluster. Assume $2\le K<N$;
zero denominators and singleton clusters require library-specific conventions.

In geophysical point processes, physical structures concentrate
along non-convex, lower-dimensional fault surfaces, curved slip
planes, and branching rupture networks. Evaluating such manifolds
with Euclidean centroid metrics leads to three distinct mathematical
pathologies:
\begin{enumerate}[noitemsep,topsep=2pt]
  \item \textbf{Non-Convex Fracture Penalisation:} Centroids
    $\mathbf{c}_k$ of some curved structures can fall outside
    the support $\mathcal{M}$. Centroid scatter may then be a poor
    proxy for density connectivity. Neither off-support centroids
    nor divergence of scatter follows from non-convexity alone.
  \item \textbf{Spurious Fragment Merging:} If two non-convex
    clusters interleave or wrap around each other, their Euclidean
    centroid distance $d(\mathbf{c}_k, \mathbf{c}_l)$ can approach
    zero even when they are separated by an absolute density void.
    Classical metrics misinterpret this as severe cluster overlap.
  \item \textbf{Absence of Statistical Scale Invariance:} None of
    the classical indices account for Poisson point-process sampling
    variance or local intrinsic dimension. Consequently, they cannot
    provide a calibrated null hypothesis or a rigorous statistical
    significance threshold.
\end{enumerate}

These limitations establish the necessity of an evidence-bounded,
mathematically rigorous framework founded directly on point-process
theory, non-parametric manifold density estimation, and topological
persistence.

% ==============================================================================
% SECTION 2: POINT PROCESS FOUNDATIONS
% ==============================================================================
\section{Measure-Theoretic and Stochastic Point Process Foundations}
\label{sec:point_process_foundations}

\subsection{Measure-Theoretic Point Process Formulation}
\label{subsec:point_process_theory}

We formalize seismic catalogs within the rigorous framework of
modern stochastic point-process theory
\parencite{daley2003,daley2008}. Let $(\Omega, \mathcal{A},
\mathbb{F}, \mathbb{P})$ be a complete filtered probability space
equipped with an internal filtration $\mathbb{F} =
(\mathcal{F}_t)_{t \in [0, T_{\mathrm{obs}}]}$, where $\mathcal{F}_t
= \sigma(\{N(A) : A \in \mathcal{B}([0, t])\})$ encodes the
historical trajectory of event arrivals up to time $t$
(\textbf{Level~2: Constitutive and Statistical Models}).

\begin{definition}[Counting Measure and Trajectory Process]
  \label{def:counting_measure}
  A temporal point process on the bounded observation interval $[0,
  T_{\mathrm{obs}}] \subset \mathbb{R}_{\ge 0}$ is a measurable
  mapping from $(\Omega, \mathcal{A})$ into the space of locally
  finite counting measures $(\mathcal{N},
  \mathcal{B}(\mathcal{N}))$. For any Borel set $A \in
  \mathcal{B}([0, T_{\mathrm{obs}}])$, the random counting measure
  is defined as:
  \begin{equation}
    N(A) \coloneqq \sum_{i=1}^{N_{\mathrm{total}}} \delta_{t_i}(A),
    \qquad N_{\mathrm{total}} \in \mathbb{N}_0,
    \label{eq:counting_measure_def}
  \end{equation}
  where $\delta_{t_i}$ is the Dirac measure concentrated at event
  occurrence time $t_i \in [0, T_{\mathrm{obs}}]$. The associated
  continuous-time counting process is $N_t \coloneqq N([0, t]) =
  \sum_{i \ge 1} \mathbf{1}_{\{t_i \le t\}}$.
\end{definition}

To ensure physical realism, we assume orderly point processes:
\begin{assumption}[Simple and Orderly Point Process]
  \label{ass:orderly}
  Assume simplicity (distinct event times almost surely) and orderliness,
  satisfying the Korolyuk condition:
  \begin{equation}
    \lim_{\Delta t \to 0^+} \frac{\mathbb{P}\bigl(N([t, t+\Delta t))
    \ge 2\bigr)}{\Delta t} = 0, \qquad \forall t \in [0, T_{\mathrm{obs}}].
    \label{eq:korolyuk_condition}
  \end{equation}
  By the separate simplicity assumption, concurrent events have
  probability zero:
  $\mathbb{P}(t_i = t_j) = 0$ for all $i \ne j$.
\end{assumption}

\begin{theorem}[Doob--Meyer Decomposition and Predictable Intensity]
  \label{thm:doob_meyer}
  Let $N_t$ be a right-continuous integrable counting process, with $N_0=0$,
  adapted to a filtration satisfying the usual completeness and
  right-continuity conditions \parencite{meyer1962}. There exists a unique
  non-decreasing predictable compensator $\Lambda_t$, with
  $\Lambda_0=0$, such that:
  \begin{equation}
    M_t \coloneqq N_t - \Lambda_t, \qquad t \in [0, T_{\mathrm{obs}}],
    \label{eq:doob_meyer}
  \end{equation}
  is a uniformly integrable $(\mathbb{F}, \mathbb{P})$-martingale.
  Orderliness alone does not imply continuity of the compensator.
  Assume here that $\Lambda_t$ is absolutely continuous with respect to the
  Lebesgue measure, there exists an $\mathbb{F}$-predictable
  non-negative process $\lambda^*(t)$ such that:
  \begin{equation}
    \Lambda_t = \int_0^t \lambda^*(s) \, \mathrm{d}s, \qquad
    \lambda^*(t) = \lim_{\Delta t \to 0^+}
    \frac{\mathbb{E}\bigl[N([t, t+\Delta t)) \mid
    \mathcal{F}_{t^-}\bigr]}{\Delta t}.
    \label{eq:compensator_intensity}
  \end{equation}
  The infinitesimal limit requires differentiation and conditional-expectation
  regularity; it is not asserted at every time point.
\end{theorem}

\subsection{The Inhomogeneous Poisson Process (IPP)}
\label{subsec:ipp_theory}

When arrival probabilities depend solely on current time $t$ rather
than past filtration history, the point process reduces to an
Inhomogeneous Poisson Process.

\begin{definition}[Inhomogeneous Poisson Process]
  \label{def:ipp}
  A point process $N(\cdot)$ is an Inhomogeneous Poisson Process
  with deterministic intensity function $\lambda: [0,
  T_{\mathrm{obs}}] \to \mathbb{R}_{\ge0}$, with finite total intensity, if:
  \begin{enumerate}[label=(\roman*),noitemsep]
    \item For any Borel set $A \in \mathcal{B}([0,
      T_{\mathrm{obs}}])$, the random count $N(A)$ follows a Poisson
      distribution with parameter $\Lambda(A) \coloneqq \int_A
      \lambda(t) \, \mathrm{d}t$:
      \begin{equation}
        \mathbb{P}\bigl(N(A) = k\bigr) = \frac{\Lambda(A)^k}{k!}
        e^{-\Lambda(A)}, \qquad k \in \mathbb{N}_0.
        \label{eq:poisson_count_dist}
      \end{equation}
    \item For any finite collection of disjoint Borel sets $A_1,
      A_2, \dots, A_m$, the random variables $N(A_1), N(A_2), \dots,
      N(A_m)$ are mutually independent.
  \end{enumerate}
\end{definition}

\begin{theorem}[Campbell's Theorem and Moment Generating Function]
  \label{thm:campbell}
  Let $N(\cdot)$ be an \ac{IPP} with intensity $\lambda(t)$ and let
  $g: [0, T_{\mathrm{obs}}] \to \mathbb{R}$ be a bounded measurable
  function \parencite{kingman1993}. The linear functional $S_g \coloneqq
  \sum_{i=1}^{N(T_{\mathrm{obs}})} g(t_i) =
  \int_0^{T_{\mathrm{obs}}} g(t) \, N(\mathrm{d}t)$ has exact
  expectation and variance given by:
  \begin{equation}
    \mathbb{E}[S_g] = \int_0^{T_{\mathrm{obs}}} g(t) \lambda(t) \,
    \mathrm{d}t, \qquad
    \Var(S_g) = \int_0^{T_{\mathrm{obs}}} g^2(t) \lambda(t) \, \mathrm{d}t.
    \label{eq:campbell_mean_var}
  \end{equation}
  Furthermore, the characteristic functional of $S_g$ evaluates to:
  \begin{equation}
    \mathbb{E}\bigl[\exp(\mathrm{i}\theta S_g)\bigr] = \exp\left(
      \int_0^{T_{\mathrm{obs}}} \bigl(e^{\mathrm{i}\theta g(t)} -
    1\bigr)\lambda(t) \, \mathrm{d}t \right), \qquad \theta \in \mathbb{R}.
    \label{eq:campbell_char_fn}
  \end{equation}
\end{theorem}

\begin{proposition}[Derivation of the Point-Process Likelihood]
  \label{prop:ipp_likelihood}
  Let $\mathcal{C} = \{t_1, t_2, \dots, t_N\}$ be an observed
  realization of ordered event arrival times on $[0,
  T_{\mathrm{obs}}]$. The likelihood functional
  $\mathcal{L}(\lambda)$ below is an ordered-event density using a fixed
  time unit, or a likelihood up to a parameter-independent reference factor:
  \begin{equation}
    \mathcal{L}(\lambda) = \left( \prod_{i=1}^N \lambda(t_i) \right)
    \exp\left( -\int_0^{T_{\mathrm{obs}}} \lambda(s) \, \mathrm{d}s \right),
    \label{eq:ipp_likelihood}
  \end{equation}
  yielding the exact continuous log-likelihood:
  \begin{equation}
    \ell(\lambda) = \ln \mathcal{L}(\lambda) = \sum_{i=1}^N \ln
    \lambda(t_i) - \int_0^{T_{\mathrm{obs}}} \lambda(s) \, \mathrm{d}s.
    \label{eq:ipp_log_likelihood}
  \end{equation}
  Relative to a Poisson reference law of positive rate $q(t)$, the exact
  log derivative is $\int\ln(\lambda/q)\,\mathrm{d}N
  -\int(\lambda-q)\,\mathrm{d}t$.
\end{proposition}

\begin{proof}
  For this bin-limit argument, assume intensity continuous and positive
  at the event times of a fixed simple finite realization.
  Partition the observation window $[0, T_{\mathrm{obs}}]$ into $M$
  disjoint subintervals $I_m = [s_{m-1}, s_m)$ of equal duration
  $\Delta = T_{\mathrm{obs}}/M$. By independent increments, the
  joint probability across the partition is:
  \begin{equation}
    \mathbb{P}\bigl(N(I_1) = k_1, \dots, N(I_M) = k_M\bigr) =
    \prod_{m=1}^M \frac{\Lambda(I_m)^{k_m}}{k_m!}
    \exp\bigl(-\Lambda(I_m)\bigr).
  \end{equation}
  As $M\to\infty$, a sufficiently fine partition of this fixed
  simple realization contains at most one event per bin. For
  intervals containing an event at $t_i \in I_m$, $\Lambda(I_m) =
  \lambda(t_i)\Delta + o(\Delta)$, contributing $\lambda(t_i)\Delta
  \exp(-\lambda(t_i)\Delta)$. For empty intervals, the factor is
  $\exp(-\Lambda(I_m))$. Factoring the product and taking the
  continuous limit as $\Delta \to 0$ cancels the Lebesgue
  discretization measure $\Delta^N$, yielding \cref{eq:ipp_likelihood}.
\end{proof}

\begin{theorem}[Maximum Entropy Characterisation of the \ac{IPP}]
  \label{thm:max_entropy_ipp}
  Let $\mathcal{P}$ be the class of all point processes on $[0,
  T_{\mathrm{obs}}]$ having fixed expected intensity measure
  $\mathbb{E}[N(A)] = \int_A \lambda(t) \, \mathrm{d}t$. Under the
  Jaynes maximum entropy principle \parencite{jaynes1957}, restrict to
  laws absolutely continuous with respect to a Poisson reference law
  $Q$ of positive rate $q(t)$, with finite relative entropies and
  integrable $\lambda|\ln(\lambda/q)|$. The \ac{IPP} maximizes
  negative relative entropy:
  \begin{equation}
    -D(P\parallel Q)\le-D(P_\lambda\parallel Q)
    =-\int_0^{T_{\mathrm{obs}}}
    [\lambda\ln(\lambda/q)-\lambda+q]\,\mathrm{d}t,
    \label{eq:max_entropy_bound}
  \end{equation}
  with equality if and only if $P=P_\lambda$. This is a
  reference-dependent entropy characterization.
\end{theorem}

\subsection{Statistical Physics and Grand Canonical Ensemble Isomorphism}
\label{subsec:grand_canonical_isomorphism}

An illuminating bridge connects stochastic point processes to
equilibrium statistical mechanics. Consider a one-dimensional gas of
non-interacting classical particles located at positions $t_i \in
[0, T_{\mathrm{obs}}]$ within a grand canonical ensemble where both
total particle number $N$ and energy fluctuate (\textbf{Level~5:
Theoretical Baselines}).

\begin{proposition}[Grand Canonical Isomorphism]
  \label{prop:grand_canonical}
  Let $\lambda(t)$ denote the Poisson intensity. Define the
  microscopic fugacity $z(t)$ and chemical potential
  $\mu_{\mathrm{chem}}(t)$ (in thermal units
  $k_{\mathrm{B}}\mathcal{T}_{\mathrm{therm}} = 1$) as:
  \begin{equation}
    z(t)=\lambda(t)/\lambda_0,\qquad
    \mu_{\mathrm{chem}}(t)=\ln z(t),\qquad\lambda_0>0.
    \label{eq:fugacity_def}
  \end{equation}
  The grand partition function $\Xi$ of the system evaluates to:
  \begin{equation}
    \Xi = \sum_{N=0}^\infty \frac{1}{N!} \left(
    \int_0^{T_{\mathrm{obs}}} \lambda_0z(t) \, \mathrm{d}t \right)^N =
    \exp\left( \int_0^{T_{\mathrm{obs}}} \lambda(t) \, \mathrm{d}t
    \right) = e^{\Lambda([0, T_{\mathrm{obs}}])}.
    \label{eq:grand_partition_function}
  \end{equation}
  The thermodynamic grand potential $\Omega_{\mathrm{grand}}$ is:
  \begin{equation}
    \Omega_{\mathrm{grand}} \coloneqq -\ln \Xi =
    -\int_0^{T_{\mathrm{obs}}} \lambda(t) \, \mathrm{d}t.
    \label{eq:grand_potential}
  \end{equation}
  Substituting \cref{eq:fugacity_def,eq:grand_potential} into the point-process
  log-likelihood in \cref{eq:ipp_log_likelihood} establishes exact equivalence:
  \begin{equation}
    \ell(\lambda)-N\ln\lambda_0 = \sum_{i=1}^N \mu_{\mathrm{chem}}(t_i) +
    \Omega_{\mathrm{grand}} \equiv -\mathcal{F}_{\mathrm{stat}}[\lambda].
    \label{eq:stat_mech_likelihood_equiv}
  \end{equation}
  Here $\lambda_0$ is a reference rate with the same units as $\lambda$.
  Maximizing likelihood minimizes the full negative log-likelihood,
  not the grand potential alone.
\end{proposition}

\subsection{Piecewise-Constant Rate Models and Optimal PELT Segmentation}
\label{subsec:pelt_segmentation}

To model non-stationary baseline fluctuations without imposing rigid
parametric shapes, we formulate the piecewise-constant Poisson
model. Let the temporal domain $[0, T_{\mathrm{obs}}]$ be
partitioned by $K$ changepoints into $K+1$ intervals:
\begin{equation}
  0 = \tau_0 < \tau_1 < \dots < \tau_K < \tau_{K+1} =
  T_{\mathrm{obs}}, \qquad I_j = (\tau_j, \tau_{j+1}].
  \label{eq:changepoint_partition}
\end{equation}
Within each interval $I_j$, the rate is constant: $\lambda(t) \equiv
\lambda_j \ge 0$. Let $n_j = N(I_j) = \sum_{i=1}^N \mathbf{1}_{\{t_i
\in I_j\}}$ denote the event count and $\Delta_j = \tau_{j+1} -
\tau_j$ the interval duration.

\begin{theorem}[Profile Log-Likelihood and Exact MLE]
  \label{thm:piecewise_mle}
  The joint log-likelihood of the changepoint partition decomposes additively:
  \begin{equation}
    \ell(\{\lambda_j\}, \{\tau_j\}) = \sum_{j=0}^K \left[ n_j \ln
    \lambda_j - \lambda_j \Delta_j \right].
    \label{eq:piecewise_loglik}
  \end{equation}
  For fixed changepoints $\{\tau_j\}$, the maximum likelihood rate
  estimators are decoupled; for $n_j>0$ the log-likelihood is strictly concave:
  \begin{equation}
    \frac{\partial \ell}{\partial \lambda_j} = \frac{n_j}{\lambda_j}
    - \Delta_j = 0 \implies \hat{\lambda}_j = \frac{n_j}{\Delta_j}.
    \label{eq:mle_lambda_j}
  \end{equation}
  For $n_j=0$, the maximizing rate is zero if allowed; with strictly
  positive rates there is only a boundary supremum. Use $0\ln0=0$.
  Substituting $\hat{\lambda}_j$ back into
  \cref{eq:piecewise_loglik} yields the exact profile log-likelihood:
  \begin{equation}
    \ell_{\mathrm{profile}}(\{\tau_j\}) = \sum_{j=0}^K \left[ n_j
    \ln\left(\frac{n_j}{\Delta_j}\right) - n_j \right].
    \label{eq:profile_loglik}
  \end{equation}
\end{theorem}

To prevent over-segmentation, the optimal partition minimizes the
penalized cost function based on the \ac{BIC}:
\begin{equation}
  \mathcal{C}_{\mathrm{BIC}}(\{\tau_j\}) = -2
  \ell_{\mathrm{profile}}(\{\tau_j\}) + (2K + 1)\ln N,
  \label{eq:bic_segmentation}
\end{equation}
where $2K+1$ counts changepoint positions and rates in this proposed
penalty convention. Changepoints are nonregular parameters, so classical
BIC guarantees do not follow from this count. A fixed candidate grid
or minimum segment duration excludes arbitrarily narrow
event-containing intervals.

\begin{definition}[Pruned Exact Linear Time (PELT) Recursion]
  \label{def:pelt}
  Define the segment cost function $\mathcal{C}(I) = -2[n_I \ln(n_I
  / |I|) - n_I]$ and let $F(s)$ denote the minimum penalized cost on
  $[0, t_s]$. By Bellman's principle of optimality:
  \begin{equation}
    F(s) = \min_{v \in \mathcal{R}_s} \left[ F(v) +
    \mathcal{C}((t_v, t_s]) + \beta \right], \qquad \beta = 2\ln N,
    \label{eq:pelt_bellman}
  \end{equation}
  where $\mathcal{R}_s \subseteq \{0, \dots, s-1\}$ is the pruned
  set of candidate changepoints. Following \citet{killick2012}, if
  there exists a constant $K_{\mathrm{prune}} \ge 0$ such that for
  all $v < s < u$:
  \begin{equation}
    \mathcal{C}((t_v, t_s]) + \mathcal{C}((t_s, t_u]) +
    K_{\mathrm{prune}} \le \mathcal{C}((t_v, t_u]),
    \label{eq:pelt_condition}
  \end{equation}
  then the candidate pruning rule:
  \begin{equation}
    \text{Prune } v \text{ if } F(v) + \mathcal{C}((t_v, t_s]) +
    K_{\mathrm{prune}} \ge F(s),
    \label{eq:pelt_pruning_rule}
  \end{equation}
  preserves the optimum value, though tied partitions can be discarded.
  Expected $\mathcal O(N)$ requires additional stochastic and cost
  conditions \parencite{killick2012}; worst-case time is quadratic.
  Any minimum-length constraints must preserve feasibility of the
  dominating alternative used in the proof.
\end{definition}

% ==============================================================================
% SECTION 3: INTRINSIC DIMENSION AND POINT-ADAPTIVE DENSITY ESTIMATION
% ==============================================================================
\section{Intrinsic Dimension and Point-Adaptive Density Estimation}
\label{sec:pak_theory}

Non-parametric density estimation in complex, high-dimensional
observational spaces faces two severe geometric challenges: the
curse of dimensionality and severe spatial heterogeneity. When
physical processes concentrate on low-dimensional
submanifolds---such as seismic hypocentres aligning along tectonic
fault planes---evaluating distances using the ambient Euclidean
dimension $D$ introduces exponential volume distortions.
Furthermore, fixing a constant neighbourhood size $k$ enforces an
intractable global bias--variance tradeoff. This section develops
the foundational geometry and statistical estimation theory of the
Point-Adaptive $k$-NN (\pak{}) framework, establishing closed-form
estimators, asymptotic optimality bounds, and rigorous uncertainty propagation.

\subsection{TWO-NN Intrinsic Dimension Estimation}
\label{subsec:two_nn_dimension}

The non-parametric estimation of intrinsic topological dimension $d$
builds upon $k$-NN dimensional scaling \parencite{levina2004,facco2017}.
Let $\mathcal{X} = \{\xvec_1, \xvec_2, \dots, \xvec_N\} \subset
\mathbb{R}^D$ be a set of $N$ independent, identically distributed
observations sampled from an underlying probability density
$\rho(\xvec)$ supported on a smooth Riemannian manifold
$\mathcal{M}$ of intrinsic topological dimension $d \le D$
(\textbf{Level~2: Idealized Manifold Geometry}).

\begin{assumption}[Local Poisson Homogeneity]
  \label{ass:local_poisson}
  For each point $\xvec_i \in \mathcal{X}$, there exists a local
  metric ball $B(\xvec_i, r)$ of radius $r \in \mathbb{R}_{>0}$
  within which the underlying probability density is locally
  constant: $\rho(\xvec) \equiv \rho_i > 0$. The generation of
  neighbouring observations within $B(\xvec_i, r)$ is modeled as a
  homogeneous spatial Poisson process with intensity $\lambda_i = N \rho_i$.
  Exact distance laws below further use boundary-free Euclidean ball
  geometry and homogeneity across the relevant neighbour distances.
  For a finite sample on a curved or truncated support these are local
  approximations, not automatic exact distributional identities.
\end{assumption}

Under \cref{ass:local_poisson}, the expected number of
observations in a $d$-dimensional ball of radius $r$ centered at
$\xvec_i$ is given by $\Lambda_i(r) = N \rho_i \omega_d r^d$, where
$\omega_d$ denotes the Lebesgue measure of the unit $d$-ball:
\begin{equation}
  \omega_d \coloneqq \frac{\pi^{d/2}}{\Gamma(d/2 + 1)}.
  \label{eq:unit_ball_volume}
\end{equation}
Let $r_{i,1}$ and $r_{i,2}$ denote the Euclidean distances from
$\xvec_i$ to its first and second nearest neighbours in $\mathcal{X}
\setminus \{\xvec_i\}$, with $0 < r_{i,1} < r_{i,2}$
(\textbf{Level~1: Raw Distance Observations}).

\begin{theorem}[Pareto Distribution of the 2NN Distance Ratio]
  \label{thm:pareto_ratio}
  Let $\mu_i \coloneqq r_{i,2} / r_{i,1} \in [1, \infty)$ denote the
  distance ratio to the first two nearest neighbours. Under
  \cref{ass:local_poisson}, the transformed variable
  $\eta_i \coloneqq \mu_i^d$ follows a standard Pareto distribution
  with scale parameter $x_m = 1$ and shape parameter $\alpha = 1$:
  \begin{equation}
    \mathbb{P}(\eta_i > t) = t^{-1}, \qquad t \ge 1.
    \label{eq:pareto_eta}
  \end{equation}
  Equivalently, $\mu_i$ is distributed according to a Pareto law
  with shape parameter $d$:
  \begin{align}
    F_\mu(\mu) &= 1 - \mu^{-d}, \qquad \mu \in [1, \infty),
    \label{eq:pareto_cdf} \\
    f_\mu(\mu) &= d \, \mu^{-(d+1)}, \qquad \mu \in [1, \infty).
    \label{eq:pareto_pdf}
  \end{align}
  In particular, the distribution of $\mu_i$ is strictly independent
  of the local density $\rho_i$, the sample size $N$, and the unit
  ball volume $\omega_d$.
\end{theorem}

\begin{proof}
  Let $\Lambda_{i,1} = N \rho_i \omega_d r_{i,1}^d$ and
  $\Lambda_{i,2} = N \rho_i \omega_d r_{i,2}^d$ denote the
  cumulative Poisson volumes enclosed by the first two nearest
  neighbours. By the properties of a homogeneous spatial Poisson
  process, the volumes $\Lambda_{i,1}$ and $\Delta \Lambda_i
  \coloneqq \Lambda_{i,2} - \Lambda_{i,1}$ represent independent
  Exponential random variables with unit rate:
  \begin{equation}
    \Lambda_{i,1} \sim \mathrm{Gamma}(1, 1), \qquad (\Lambda_{i,2} -
    \Lambda_{i,1}) \sim \mathrm{Gamma}(1, 1).
  \end{equation}
  The joint probability density function of $(\Lambda_{i,1},
  \Lambda_{i,2})$ over the domain $0 < \lambda_1 < \lambda_2 <
  \infty$ is $p(\lambda_1, \lambda_2) = e^{-\lambda_1}
  e^{-(\lambda_2 - \lambda_1)} = e^{-\lambda_2}$. Consider the
  change of variables $(\Lambda_{i,1}, \Lambda_{i,2}) \mapsto
  (\Lambda_{i,1}, \eta_i)$, where $\eta_i =
  \Lambda_{i,2}/\Lambda_{i,1} = (r_{i,2}/r_{i,1})^d = \mu_i^d$.
  Substituting $\lambda_2 = \eta \lambda_1$, the Jacobian
  determinant is $J = \lambda_1$. The joint density of
  $(\Lambda_{i,1}, \eta_i)$ on $\lambda_1 \in (0, \infty)$ and $\eta
  \in [1, \infty)$ is $p(\lambda_1, \eta) = \lambda_1 e^{-\eta
  \lambda_1}$. Marginalizing over $\lambda_1$:
  \begin{equation}
    p(\eta) = \int_0^\infty \lambda_1 e^{-\eta \lambda_1} \,
    \mathrm{d}\lambda_1 = \frac{1}{\eta^2}, \qquad \eta \ge 1.
  \end{equation}
  Applying the transformation $\eta = \mu^d$ with derivative
  $\mathrm{d}\eta/\mathrm{d}\mu = d \, \mu^{d-1}$:
  \begin{equation}
    f_\mu(\mu) = p(\eta) \left| \frac{\mathrm{d}\eta}{\mathrm{d}\mu}
    \right| = \frac{1}{(\mu^d)^2} \cdot d \, \mu^{d-1} = d \,
    \mu^{-(d+1)}, \qquad \mu \ge 1.
  \end{equation}
  Integrating $f_\mu(\mu)$ from $1$ to $\mu$ yields $F_\mu(\mu) = 1
  - \mu^{-d}$. The common factor $N \rho_i \omega_d$ cancels
  identically in the ratio $\Lambda_{i,2}/\Lambda_{i,1}$, proving
  that $F_\mu(\mu)$ is an exact pivotal quantity for $d$ \parencite{facco2017}.
\end{proof}

\begin{proposition}[Maximum Likelihood Estimator and Asymptotic
  Variance of $d$]
  \label{prop:2nn_mle}
  Given $N$ independent distance ratios $\{\mu_1, \mu_2, \dots,
  \mu_N\}$, the Maximum Likelihood Estimator (\ac{MLE}) of intrinsic
  dimension $d$ (\textbf{Level~4: Model Estimate}) is:
  \begin{equation}
    \hat{d}_{\mathrm{MLE}} = \frac{N}{\sum_{i=1}^N \ln \mu_i}.
    \label{eq:2nn_mle}
  \end{equation}
  The Fisher information for the Pareto shape parameter is
  $\mathcal{I}(d) = 1/d^2$, and the estimator achieves the
  Cram\'er--Rao lower bound asymptotically:
  \begin{equation}
    \sqrt{N} \bigl(\hat{d}_{\mathrm{MLE}} - d\bigr)
    \xrightarrow{\mathcal{D}} \mathcal{N}\bigl(0, \, d^2\bigr),
    \qquad \Var\bigl(\hat{d}_{\mathrm{MLE}}\bigr) = \frac{d^2}{N} +
    \mathcal{O}\bigl(N^{-2}\bigr).
    \label{eq:2nn_asymptotic_var}
  \end{equation}
  These limits concern independent ideal ratios. At finite $N>1$,
  $\mathbb{E}[\hat d_{\mathrm{MLE}}]=Nd/(N-1)$; an unbiased version
  is $(N-1)/\sum_i\ln\mu_i$. Shared sample neighbourhoods generate
  dependent ratios, requiring a separate variance analysis. The source
  uses trimmed empirical-CDF regression, not this independent-ratio MLE.
\end{proposition}

\begin{theorem}[Excess Shannon Entropy and Scale Invariance]
  \label{thm:excess_entropy}
  Under the ideal fixed-rank experiment, the differential entropy of $R_k$
  contains an explicit dependence on the local density $\rho_i$:
  \begin{equation}
    H(R_k) = k + \ln \Gamma(k) - (k - 1)\psi(k) - \ln d - \frac{d -
    1}{d}\psi(k) - \frac{1}{d}\ln(N \rho_i \omega_d),
    \label{eq:shannon_rk}
  \end{equation}
  where $\psi(k) = \Gamma'(k)/\Gamma(k)$ is the digamma function. In
  contrast, the differential entropy of the 2NN ratio $\mu =
  R_2/R_1$ is strictly density-free:
  \begin{equation}
    H(\mu) \coloneqq -\int_1^\infty f_\mu(x) \ln f_\mu(x) \,
    \mathrm{d}x = 1 + \frac{1}{d} - \ln d.
    \label{eq:excess_entropy}
  \end{equation}
  Distance units are fixed. This ratio law is invariant to uniform
  rescaling, not arbitrary anisotropic or nonlinear geometric transforms.
\end{theorem}

\subsection{Adaptive Neighbourhood Selection via Sequential
Likelihood-Ratio Testing}
\label{subsec:adaptive_kstar}

Having estimated intrinsic dimension $d$, we determine for each
point $\xvec_i$ a neighbourhood selected by an operational stopping
statistic. It does not certify local homogeneity (\cref{ass:local_poisson}).

Let $r_{i,1} < r_{i,2} < \dots < r_{i,k}$ denote ordered
nearest-neighbour Euclidean distances. Concentric shell volumes are defined as:
\begin{equation}
  w_{i,j} \coloneqq \omega_d \bigl(r_{i,j}^d - r_{i,j-1}^d\bigr),
  \qquad j \in \{1, 2, \dots, k\}, \quad (r_{i,0} \equiv 0).
  \label{eq:shell_volume_def}
\end{equation}
Let $v_{i,j} \coloneqq \sum_{m=1}^j w_{i,m} = \omega_d r_{i,j}^d$ be
the volume of the inner sphere containing $j$ neighbours, and
$v_{i,k} - v_{i,j}$ the volume of the outer shell containing the
remaining $k - j$ neighbours.

\begin{theorem}[Log-Likelihood Ratio Test Statistic and Wilks' Limiting Law]
  \label{thm:wilks_deviance}
  Under the null hypothesis $H_0$ (homogeneous Poisson process with
  rate $\rho_0$ throughout $B(\xvec_i, r_{i,k})$), the maximized
  log-likelihood is $\ell_0 = k \ln(k/v_{i,k}) - k$. Under the
  alternative $H_1$ (density changepoint at shell $j$ with rates
  $\rho_1 \ne \rho_2$), the maximized log-likelihood is $\ell_1 = j
  \ln(j/v_{i,j}) + (k-j)\ln((k-j)/(v_{i,k}-v_{i,j})) - k$. The
  log-likelihood ratio deviance statistic is:
  \begin{equation}
    \Lambda_{k,j} \coloneqq 2(\ell_1 - \ell_0) = 2 \left[ j
      \ln\left( \frac{\hat{\rho}_1}{\hat{\rho}_0} \right) + (k - j)
    \ln\left( \frac{\hat{\rho}_2}{\hat{\rho}_0} \right) \right],
    \label{eq:deviance_statistic}
  \end{equation}
  where $\hat{\rho}_0 = k/v_{i,k}$, $\hat{\rho}_1 = j/v_{i,j}$, and
  $\hat{\rho}_2 = (k-j)/(v_{i,k} - v_{i,j})$. A Wilks reference
  \parencite{wilks1938} requires a prespecified regular two-rate
  experiment with positive rates and both sample sizes increasing:
  \begin{equation}
    \Lambda_{k,j} \xrightarrow{\mathcal{D}} \chi^2_1, \qquad
    \text{as } k \to \infty.
    \label{eq:wilks_convergence}
  \end{equation}
  Searching over splits or ranks does not preserve this reference law
  without additional calibration. This two-region deviance is a
  theoretical baseline, not the operational stopping statistic below.
\end{theorem}

\begin{definition}[Sequential Halting Rule for Adaptive $\kstar(i)$]
  \label{def:kstar_stopping}
  In the inspected scoring implementation, testing starts at
  $k_{\min}=4$ and considers ranks $k<k_{\max}$ using
  \begin{equation}
    D_{i,k}=2(k\hat d+1)\ln(r_{i,k+1}/r_{i,k})
    -2\ln\hat d-2\ln k.
  \end{equation}
  It does not maximize \cref{eq:deviance_statistic} over split indices.
  The selected neighbourhood is:
  \begin{equation}
    \kstar(i)=\min\{k_{\min}\le k<k_{\max}:D_{i,k}>D_{\mathrm{thr}}\},
    \label{eq:kstar_stopping_rule}
  \end{equation}
  with $\kstar(i) = k_{\max}$ if no statistic exceeds threshold. The
  critical threshold is set to $D_{\mathrm{thr}} \coloneqq \num{23.93}$,
  corresponding to:
  \begin{equation}
    \mathbb{P}\bigl(\chi^2_1 > \num{23.928}\bigr) = \num{1.0e-6}.
    \label{eq:dthr_pvalue}
  \end{equation}
  The $\chi^2_1$ tail is only a reference probability, not demonstrated
  calibration of $D_{i,k}$ or its sequential search. A Bonferroni
  \ac{FWER} guarantee requires valid individual tests and a specified
  comparison family; it is not established here.
\end{definition}

\subsection{Point-Adaptive Maximum Likelihood Density Estimation}
\label{subsec:pak_density_estimation}

Given $\kstar(i)$, the \pak{} estimator evaluates the complete
configuration of all $\kstar(i)$ nearest neighbours \parencite{derrico2021}.

\begin{theorem}[Poisson Shell Likelihood and Closed-Form MLE]
  \label{thm:shell_mle}
  Fix $k=\kstar(i)$ independently of the observations, with known
  dimension and homogeneous boundary-free geometry. Let
  $S_i=\omega_dr_{i,k}^d$ be the rank-$k$ stopping volume.
  The joint density of ordered arrival volumes
  $0<v_{i,1}<\cdots<v_{i,k}=S_i$, relative to their Lebesgue volume
  elements, is:
  \begin{equation}
    \mathcal{L}(\rho_i)=(N\rho_i)^k\exp(-N\rho_iS_i).
    \label{eq:joint_shell_likelihood}
  \end{equation}
  Reparameterizing in log-intensity $u_i \coloneqq \ln(N \rho_i)$
  yields the log-likelihood:
  \begin{equation}
    \ell(u_i)=ku_i-e^{u_i}S_i.
    \label{eq:log_lik_ui}
  \end{equation}
  The score and Hessian are:
  \begin{equation}
    \ell'(u_i) = \kstar(i) - e^{u_i} S_i, \qquad \ell''(u_i) =
    -e^{u_i} S_i < 0.
    \label{eq:score_hessian_u}
  \end{equation}
  The unique maximum likelihood estimators in intensity and
  probability density space are:
  \begin{equation}
    e^{\hat{u}_i} \coloneqq \frac{\kstar(i)}{S_i} \iff \hat{\rho}_{i,
    \mathrm{MLE}} \coloneqq \frac{\kstar(i)}{N \omega_d r_{i, \kstar(i)}^d}.
    \label{eq:mle_densities}
  \end{equation}
  The source's nominal root is log intensity $\ln(k/S_i)$, without
  probability-density normalization by $N$. Its use of adaptive $k^*$
  does not by itself justify the fixed-rank sampling guarantees below.
\end{theorem}

\begin{proposition}[Bias Correction and UMVUE Density Estimator]
  \label{prop:umvue_density}
  Under this fixed-rank experiment, $N \rho_i S_i \sim
  \mathrm{Gamma}(\kstar(i), 1)$, so the raw
  MLE has positive bias: $\mathbb{E}[\hat{\rho}_{i, \mathrm{MLE}}] =
  \frac{\kstar(i)}{\kstar(i) - 1} \rho_i$. The unique \ac{UMVUE} is:
  \begin{equation}
    \hat{\rho}_{i, \mathrm{UB}} \coloneqq \frac{\kstar(i) - 1}{N \omega_d
    r_{i, \kstar(i)}^d}, \qquad \Var\bigl(\hat{\rho}_{i,
    \mathrm{UB}}\bigr) = \frac{\rho_i^2}{\kstar(i) - 2}, \quad
    \text{for } \kstar(i) \ge 3.
    \label{eq:umvue_density}
  \end{equation}
  The Fisher information is $\mathcal{I}(\rho_i) =
  \kstar(i)/\rho_i^2$. The estimator $\hat{\rho}_{i, \mathrm{UB}}$
  asymptotically saturates the Cram\'er--Rao lower bound:
  \begin{equation}
    \operatorname{Eff}\bigl(\kstar(i)\bigr) \coloneqq
    \frac{\mathrm{CRLB}}{\Var(\hat{\rho}_{i, \mathrm{UB}})} =
    \frac{\kstar(i) - 2}{\kstar(i)} \xrightarrow{\kstar \to \infty} 1.
    \label{eq:asymp_efficiency}
  \end{equation}
  Completeness and sufficiency of the fixed-rank Gamma stopping volume
  justify the UMVUE statement. It does not automatically survive
  data-dependent $k^*$ or estimated dimension; the source does not
  implement this bias correction.
\end{proposition}

\subsection{Uncertainty Propagation and the Conservative Error Model}
\label{subsec:uncertainty_propagation}

For the ideal fixed-rank unbiased estimator,
$\Var(\hat{\rho}_i) = \rho_i^2/(\kstar(i) - 2)$ scales
quadratically with unknown density $\rho_i$, raw density contrasts
are severely heteroscedastic.

\begin{proposition}[Variance-Stabilizing Logarithmic Transformation
  via Delta Method]
  \label{prop:delta_method_log}
  Let $\phi_i \coloneqq \ln \hat{\rho}_{i, \mathrm{UB}}$ denote the
  log-density. Applying the Delta method \parencite{cramer1946} with
  $g(\rho) = \ln \rho$:
  \begin{equation}
    \Var\bigl(\ln \hat{\rho}_{i, \mathrm{UB}}\bigr) \approx \bigl[
    g'(\rho_i) \bigr]^2 \Var\bigl(\hat{\rho}_{i, \mathrm{UB}}\bigr)
    = \left( \frac{1}{\rho_i} \right)^{\!2}
    \frac{\rho_i^2}{\kstar(i) - 2} = \frac{1}{\kstar(i) - 2}.
    \label{eq:delta_method_result}
  \end{equation}
  This is a first-order fixed-rank approximation, not an exact adaptive
  uncertainty identity. The exact fixed-rank log variance is given next.
\end{proposition}

\begin{theorem}[Exact Gamma Variance via Trigamma Function]
  \label{thm:trigamma_variance}
  From $\hat{u}_i = \ln \kstar(i) - \ln S_i$ with $S_i = \Lambda_i /
  (N \rho_i)$ and $\Lambda_i \sim \mathrm{Gamma}(\kstar(i), 1)$, the
  exact variance for a fixed rank is:
  \begin{equation}
    \Var[\hat{u}_i] = \Var[\ln \Lambda_i] =
    \psi_1\bigl(\kstar(i)\bigr) = \frac{1}{\kstar(i)} +
    \frac{1}{2(\kstar(i))^2} + \mathcal{O}\bigl((\kstar(i))^{-3}\bigr),
    \label{eq:trigamma_exact}
  \end{equation}
  where $\psi_1(z) = \frac{\mathrm{d}^2}{\mathrm{d}z^2}\ln
  \Gamma(z)$ is the trigamma function.
\end{theorem}

\begin{definition}[Conservative PAk Log-Density Error Formulation]
  \label{def:pak_error_formula}
  In the computational implementation of \pak{}, the per-point
  standard error is evaluated via the conservative model:
  \begin{equation}
    \hat{\sigma}_i \coloneqq \sqrt{\frac{4\kstar(i) +
    2}{\kstar(i)\bigl(\kstar(i) - 1\bigr)}} =
    \sqrt{\frac{4}{\kstar(i)} +
    \mathcal{O}\bigl((\kstar(i))^{-2}\bigr)} \approx
    \frac{2}{\sqrt{\kstar(i)}}.
    \label{eq:pak_conservative_error}
  \end{equation}
\end{definition}

\begin{remark}[Epistemic Justification of the \num{4}$\times$ Variance Margin]
  \label{rem:conservative_justification}
  As demonstrated in \Cref{tab:variance_comparison}, the
  operational error model in \cref{eq:pak_conservative_error} exceeds
  the nominal CRLB ($1/\kstar$) by a factor of $\approx \num{4}$ in
  variance (factor of $\approx \num{2}$ in standard error). Relevant
  departures from the ideal fixed-rank model include:
  (i) random fluctuations of empirical shell volumes $S_i$;
  (ii) propagation of intrinsic dimension uncertainty $\Var(\hat{d})
  \approx \hat{d}^2/N$ through volume exponentiation;
  (iii) adaptive selection bias from sequential LRT search; and
  (iv) residual continuous density gradients violating strict local
  homogeneity.
  The formula alone does not bound these departures or prove coverage,
  Type~I error control, or calibration of selected density saddles.
\end{remark}

\begin{table}[htbp]
  \centering
  \small
  \begin{threeparttable}
    \caption{\textbf{Comparison of Log-Density Variance Estimators
      Across Neighbourhood Sizes}. Target/Scope: Evaluation of
      analytical variance formulas for log-density
      estimator $\ln \hat{\rho}_i$ across characteristic sizes $\kstar
      \in \{4, 10, 50, 100\}$. Key Phenomenon: The conservative \pak{}
      formulation converges asymptotically to $4\times$ the
      Cram\'er--Rao lower bound; the standard-error ratio tends to two,
      but is not exactly two at finite rank. These are formula
    evaluations for a fixed-rank Poisson baseline, not coverage results.}
    \label{tab:variance_comparison}
    \begin{tabular*}{\linewidth}{@{\extracolsep{\fill}}l
        S[table-format=1.4] S[table-format=1.4] S[table-format=1.4]
      S[table-format=1.4] S[table-format=1.2]@{}}
      \toprule
      {Rank} & {CRLB} & {Trigamma} & {Delta} &
      {\pak{}} & {Ratio} \\
      {$k$} & {$1/k$} & {$\psi_1(k)$} & {$1/(k-2)$} &
      {$(4k+2)/[k(k-1)]$} & {$\hat\sigma^2/\mathrm{CRLB}$} \\
      \midrule
      \num{4}   & 0.2500 & 0.2838 & 0.5000 & 1.5000 & 6.00 \\
      \num{10}  & 0.1000 & 0.1052 & 0.1250 & 0.4667 & 4.67 \\
      \num{50}  & 0.0200 & 0.0202 & 0.0208 & 0.0824 & 4.12 \\
      \num{100} & 0.0100 & 0.0101 & 0.0102 & 0.0406 & 4.06 \\
      \bottomrule
    \end{tabular*}
    \begin{tablenotes}[flushleft]
      \footnotesize
    \item Target \& Scope: [Derived Metric] Formula evaluations at
      $k\in\{4,10,50,100\}$, with fixed $k$ for the theoretical laws;
      the operational formula is evaluated at $k=k^*$.
    \item Notice: As $\kstar \to \infty$, the conservative PAk error
      model converges asymptotically to four times the Cram\'er--Rao
      lower bound ($4/\kstar$). This algebraic comparison is not a
      guaranteed uncertainty margin for empirical sampling fluctuations.
    \item Evidence Anchor: Values are evaluations of the displayed
      fixed-rank formulas, not empirical variance estimates. No Poisson
      simulation coverage was verified in this review.
    \end{tablenotes}
  \end{threeparttable}
\end{table}

% ==============================================================================
% SECTION 4: SELF-EXCITING POINT PROCESSES AND SEISMOTECTONIC TRIGGERING
% ==============================================================================
\section{Self-Exciting Point Processes and Seismotectonic Triggering}
\label{sec:clustering_hawkes_etas}

While non-parametric estimators capture macroscopic rate
topographies, physical seismic catalogs exhibit microscopic
cascading: each rupture redistributes crustal stress, elevating the
nucleation probability of secondary events on adjacent faults. To
incorporate endogenous triggering, this section establishes the
mathematical foundations of marked Hawkes processes
\parencite{hawkes1971} and the \ac{ETAS} formulation
\parencite{ogata1988,ogata1998}, derives the analytical stability
criterion, and formalises the mechanical contrast between
stress-transferred aftershocks and fluid-driven swarms.

\subsection{Marked Hawkes Processes and Spatio-Temporal ETAS Formulation}
\label{subsec:clustering_marked_etas}

Let $\mathcal{S} \subset \mathbb{R}^2$ denote the planar
seismotectonic domain, $\mathcal{T} = [0, T_{\mathrm{obs}}]$ the
temporal observation horizon, and $\mathcal{M} = [m_0, \infty)$ the
magnitude mark space bounded below by catalogue completeness
threshold $m_0 \in \mathbb{R}_{>0}$ (\textbf{Level~2: Ground Truth
Completeness}). An earthquake catalogue is an ordered sequence of
marked events:
\begin{equation}
  \mathcal{C} = \bigl\{ (t_i, \mathbf{x}_i, m_i) \bigr\}_{i=1}^{N},
  \qquad 0 \le t_1 < \dots < t_N \le T_{\mathrm{obs}}, \quad
  \mathbf{x}_i \in \mathcal{S}, \quad m_i \in \mathcal{M}.
  \label{eq:clustering_catalog_def}
\end{equation}

\begin{definition}[Marked Spatio-Temporal Conditional Intensity]
  \label{def:clustering_conditional_intensity}
  Under mark separability \parencite{ogata1988,ogata1998}, magnitude
  marks are conditionally independent of history given origin time
  and epicentre:
  \begin{equation}
    \lambda(t, \mathbf{x}, m \mid \mathcal{H}_t) = \lambda(t,
    \mathbf{x} \mid \mathcal{H}_t) \, p(m),
    \label{eq:clustering_intensity_separable}
  \end{equation}
  where $\mathcal{H}_t = \{(t_j, \mathbf{x}_j, m_j) : t_j < t\}$ is
  the filtration history. In spatio-temporal \ac{ETAS}, ground
  intensity decomposes linearly:
  \begin{equation}
    \lambda(t, \mathbf{x} \mid \mathcal{H}_t) = \mu_0(\mathbf{x}) +
    \sum_{j:\, t_j < t} g(t - t_j, m_j) \, f(\mathbf{x} -
    \mathbf{x}_j \mid m_j),
    \label{eq:clustering_etas_spatiotemporal}
  \end{equation}
  where $\mu_0(\mathbf{x}) \ge 0$ is time-invariant background rate
  density, $g(\Delta t, m)$ is the temporal triggering kernel, and
  $f(\mathbf{r} \mid m)$ is spatial redistribution normalized such
  that $\int_{\mathbb{R}^2} f(\mathbf{r} \mid m) \,
  \mathrm{d}\mathbf{r} = 1$. Integrating spatial dimensions over
  all space recovers purely temporal marked \ac{ETAS}; for bounded
  $\mathcal S$, each parent contribution also carries its kernel
  retention factor $\int_{\mathcal S}f(\mathbf x-\mathbf x_j\mid m_j)
  \,\mathrm{d}\mathbf x$:
  \begin{equation}
    \lambda(t \mid \mathcal{H}_t) = \mu_0(t) + \sum_{j:\, t_j < t}
    g(t - t_j, m_j), \qquad t \in \mathcal{T}.
    \label{eq:clustering_etas_temporal}
  \end{equation}
\end{definition}

\subsection{Constitutive Kernel Functions and Mark Statistics}
\label{subsec:clustering_kernels}

Earthquake magnitudes follow the Gutenberg--Richter exponential distribution:
\begin{equation}
  p(m) = \beta \, e^{-\beta(m - m_0)}, \qquad m \ge m_0, \quad \beta
  = b \ln(10) \in \mathbb{R}_{>0}.
  \label{eq:clustering_gr_pdf}
\end{equation}
The temporal triggering kernel $g(\tau, m)$ combines the Omori--Utsu
power-law relaxation with exponential magnitude efficiency:
\begin{equation}
  g(\tau, m) = K \, (\tau + c)^{-p} \, e^{\alpha(m - m_0)}, \qquad
  \tau > 0, \quad m \ge m_0,
  \label{eq:clustering_omori_kernel}
\end{equation}
where $K \in \mathbb{R}_{>0}$ is productivity scale, $c \in
\mathbb{R}_{>0}$ regularisation time cutoff, $p \in (1, \infty)$
Omori decay exponent, and $\alpha \in (0, \beta)$ magnitude
efficiency parameter.

\subsection{Analytical Derivation of the Branching Ratio and Stability}
\label{subsec:clustering_branching_ratio}

In marked self-exciting point processes \parencite{helmstetter2002}, the
branching ratio controls whether cascades remain contained or trigger
runaway supercritical seismicity.

\begin{theorem}[Branching Ratio Formulation]
  \label{thm:clustering_branching_ratio}
  Let temporal kernel $g(\tau, m)$ follow
  \cref{eq:clustering_omori_kernel} with $p > 1$, and
  magnitudes follow \cref{eq:clustering_gr_pdf} with $\beta >
  \alpha$. The branching ratio $n_{\mathrm{b}} \in \mathbb{R}_{>0}$
  (expected first-generation offspring per event) is:
  \begin{equation}
    n_{\mathrm{b}} \coloneqq \mathbb{E}_{M}\left[ \int_0^\infty
    g(\tau, M) \, \mathrm{d}\tau \right] = \frac{K \, c^{1-p}}{p -
    1} \cdot \frac{\beta}{\beta - \alpha}.
    \label{eq:clustering_branching_ratio}
  \end{equation}
\end{theorem}

\begin{proof}
  Integrating $g(\tau, m)$ over $\tau \in [0, \infty)$ using
  substitution $u = \tau + c$:
  \begin{equation}
    \kappa(m) \coloneqq \int_0^\infty g(\tau, m) \, \mathrm{d}\tau =
    K e^{\alpha(m - m_0)} \int_c^\infty u^{-p} \, \mathrm{d}u =
    \frac{K c^{1-p}}{p - 1} e^{\alpha(m - m_0)}, \quad (p > 1).
  \end{equation}
  Integrating $\kappa(m)$ against $p(m) = \beta e^{-\beta(m - m_0)}$
  with substitution $v = m - m_0$:
  \begin{equation}
    n_{\mathrm{b}} = \frac{K c^{1-p}\beta}{p - 1}\int_0^\infty
    e^{-(\beta - \alpha)v} \, \mathrm{d}v = \frac{K c^{1-p}\beta}{(p
    - 1)(\beta - \alpha)}, \quad (\beta > \alpha),
  \end{equation}
  which establishes \cref{eq:clustering_branching_ratio}.
\end{proof}

\begin{corollary}[Critical Stability Boundaries]
  \label{cor:clustering_criticality}
  The parameter space partitions into three asymptotic regimes:
  \begin{enumerate}[label=(\roman*),noitemsep]
    \item \textbf{Subcritical ($n_{\mathrm{b}} < 1$):} Cascades
      extinguish almost surely. A stationary construction with constant
      immigration on the whole time axis has finite mean rate.
      Expected total progeny seeded by an immigrant event
      is $\mathbb{E}[Z] = \sum_{k=0}^\infty n_{\mathrm{b}}^k = (1 -
      n_{\mathrm{b}})^{-1}$. The mean stationary seismicity rate is
      $\bar{\lambda} = \mu_0 / (1 - n_{\mathrm{b}})$.
    \item \textbf{Critical ($n_{\mathrm{b}} = 1$):} Critical
      productivity $K_{\mathrm{crit}} = \frac{(p - 1)(\beta -
      \alpha)}{\beta c^{1-p}}$. Extinction is almost sure for the
      nondegenerate offspring law, while $\mathbb E[Z]=\infty$.
      A $z^{-1/2}$ size tail additionally requires finite offspring
      variance; for these unbounded marks, $\beta>2\alpha$.
    \item \textbf{Supercritical ($n_{\mathrm{b}}>1$):} Infinite
      genealogical cascades have positive probability. This alone
      does not imply infinite counts at every finite time or
      establish finite-time explosion.
  \end{enumerate}
\end{corollary}

\subsection{Physical Triggering Mechanics: Coulomb Stress, Friction,
and Diffusion}
\label{subsec:clustering_physical_mechanisms}

\paragraph{Coulomb Failure Stress Transfer ($\Delta\mathrm{CFF}$).}
Coseismic slip creates an instantaneous elastostatic stress
perturbation $\Delta\bm{\sigma}$ in the surrounding crust
\parencite{okada1992,king1994,stein1999,harris1998}. On a receiver
fault patch defined by
unit normal $\mathbf{n}$ and slip vector $\mathbf{s}$:
\begin{equation}
  \Delta\mathrm{CFF} \coloneqq \Delta\tau_{\mathrm{s}} + \mu_{\mathrm{f}}
  (\Delta\sigma_{\mathrm{n}} + \Delta p) \equiv \Delta\tau_{\mathrm{s}} +
  \mu_{\mathrm{eff}} \, \Delta\sigma_{\mathrm{n}},
  \label{eq:clustering_coulomb_cff}
\end{equation}
where $\Delta\tau_{\mathrm{s}}$ is shear stress change along slip
direction, $\Delta\sigma_{\mathrm{n}}$ normal stress change (tensile
positive), $\Delta p$ pore pressure change, $\mu_{\mathrm{f}} \in
[\num{0.4}, \num{0.85}]$ coefficient of internal friction, and
$\mu_{\mathrm{eff}} \coloneqq \mu_{\mathrm{f}}(1 - B)$ apparent friction
with Skempton coefficient $B \in [\num{0}, \num{1}]$. Nucleation is promoted
when $\Delta\mathrm{CFF} > 0$ and suppressed in stress shadows
($\Delta\mathrm{CFF} < 0$).

\paragraph{Rate-and-State Friction Derivation of Omori Decay.}
Under the rate-and-state friction formulation
\parencite{dieterich1994,ruina1983}, seismicity rate $R(t)$ under steady
stressing rate $\dot{\tau}_{\mathrm{r}}$ following stress step
$\Delta\tau = \Delta\mathrm{CFF}$ satisfies:
\begin{equation}
  \frac{\mathrm{d}R}{\mathrm{d}t} = \frac{R}{t_{\mathrm{a}}} \left(
  \frac{\dot{\tau}}{\dot{\tau}_{\mathrm{r}}} - \frac{R}{r} \right)
  \implies R(t) = \frac{r}{\left[
      \exp\left(-\frac{\Delta\tau}{A\sigma_{\mathrm{n}}}\right) - 1
  \right] \exp\left(-\frac{t}{t_{\mathrm{a}}}\right) + 1},
  \label{eq:clustering_dieterich_exact}
\end{equation}
where $t_{\mathrm{a}} \coloneqq
A\sigma_{\mathrm{n}}/\dot{\tau}_{\mathrm{r}}$ is aftershock
relaxation duration, with $A\sigma_{\mathrm n}>0$ the effective frictional
stress scale. For $t\ll t_{\mathrm a}$ and
$\Delta\tau\gg A\sigma_{\mathrm n}$, this approximates
Omori's law with $p = \num{1}$:
\begin{equation}
  R(t) \approx \frac{r \,
  \exp\left(\frac{\Delta\tau}{A\sigma_{\mathrm{n}}}\right)}{1 +
  \frac{t}{c_{\mathrm{RS}}}}, \qquad c_{\mathrm{RS}} \coloneqq
  t_{\mathrm{a}} \exp\left(-\frac{\Delta\tau}{A\sigma_{\mathrm{n}}}\right).
  \label{eq:clustering_dieterich_omori}
\end{equation}
The initial rate is exponential in stress, not necessarily the integrated
count. A magnitude-productivity law additionally requires a stress--magnitude
relation and an affected-volume model.

\paragraph{Pore-Pressure Diffusion in Fluid-Driven Swarms.}
In fluid-saturated porous elastic media, Darcy flux $\mathbf{q}
\coloneqq -\frac{k}{\eta}\nabla p \equiv -\frac{k}{\eta}\operatorname{grad} p$
and fluid mass conservation $\nabla \cdot \mathbf{q} +
S_{\mathrm{s}}\frac{\partial p}{\partial t} \equiv \operatorname{div}\mathbf{q}
+ S_{\mathrm{s}}\frac{\partial p}{\partial t} = Q(\mathbf{x}, t)$ yield the
parabolic diffusion equation for pore pressure $p(\mathbf{x}, t)$
\parencite{rice1976}:
\begin{equation}
  \frac{\partial p}{\partial t} = D \, \nabla^2 p +
  \frac{Q(\mathbf{x}, t)}{S_{\mathrm{s}}} \equiv D
  \operatorname{div}(\operatorname{grad} p)
  + \frac{Q(\mathbf{x}, t)}{S_{\mathrm{s}}}, \qquad D \coloneqq
  \frac{k}{\eta S_{\mathrm{s}}} \in \mathbb{R}_{>0},
  \label{eq:clustering_diffusion_pde}
\end{equation}
where $D$ is hydraulic diffusivity, $k$ permeability, $\eta$ dynamic
viscosity, and $S_{\mathrm{s}}$ pressure storage in inverse pressure
units; hydraulic-head specific storage requires a density/gravity
conversion. A three-dimensional impulse gives a peak-time locus
$r=\sqrt{6Dt}$ and a distinct logarithmically modified pressure contour.
The separately motivated Shapiro envelope convention
\parencite{shapiro1997} is:
\begin{equation}
  r(t) \coloneqq \sqrt{4\pi D t}.
  \label{eq:clustering_shapiro_front}
\end{equation}
A linear envelope in distance--$\sqrt t$ space is compatible with diffusion,
not unique evidence of fluid causation. The three constructions are shown
in \cref{fig:supp_diffusion_constructions}; physical interpretations appear in
(\cref{tab:clustering_aftershock_vs_swarm}).

\begin{table}[htbp]
  \centering
  \small
  \begin{threeparttable}
    \caption{\textbf{Diagnostic Seismotectonic Classification
      Matrix}. Systematic comparison between stress-triggered
      aftershock cascades and fluid-driven seismic swarms across
      physical, statistical, and kinematic criteria [Target: Cluster
        Classification; Key Phenomenon: Mechanics vs Point-Process
    Signatures; Conditions: Completeness $m \ge m_0$, Tectonic Fault Systems].}
    \label{tab:clustering_aftershock_vs_swarm}
    \begin{tabularx}{\textwidth}{l >{\raggedright\arraybackslash}X
      >{\raggedright\arraybackslash}X}
      \toprule
      \textbf{Diagnostic Criterion} & \textbf{Stress-Triggered
      Aftershocks} & \textbf{Fluid-Driven Seismic Swarms} \\
      \midrule
      \textbf{Primary Driver} & Co-seismic elastostatic stress step
      ($\Delta\mathrm{CFF}$) & Pore-pressure diffusion ($\partial p
        / \partial t = D \nabla^2 p \equiv D
      \operatorname{div}(\operatorname{grad} p)$) or aseismic slip \\
      \textbf{Energy Partitioning} & Dominant mainshock; conforms to
      B\r{a}th's law ($\Delta m \approx \num{1.2}$) & No dominant
      mainshock; multiple events of comparable magnitude \\
      \textbf{Temporal Rate Decay} & Monotonic power-law Omori
      relaxation: $\sim (\tau + c)^{-p}$ ($p \in
      \numrange{1.0}{1.3}$) & Transient asymmetric envelope (log-normal
      pulse); gradual onset \\
      \textbf{Branching Ratio ($n_{\mathrm{b}}$)} & Must be estimated
      under a specified kernel & Not determined by a swarm label alone \\
      \textbf{Gutenberg--Richter $b$-value} & Standard crustal
      baseline: $b \in \numrange{0.8}{1.1}$ & Elevated: $b \in
      \numrange{1.3}{2.5}$ due to high pore-fluid pressure \\
      \textbf{Spatial Kinematics} & Static spatial footprint bounded
      by mainshock fault rupture & Diffusive parabolic migration
      front: $r(t) \propto \sqrt{4\pi D t}$ \\
      \textbf{Optimal Model Representation} & Pure parametric
      \ac{ETAS} ($d = \num{5}$: $K, c, p, \alpha, \mu$) & Multi-scale
      Unified Model: log-normal envelope $\phi_k(t)$ + \ac{ETAS} aftershocks \\
      \textbf{Epistemic Class} & Level~5: Elastostatic
      rate-and-state mechanics & Level~5: Poroelastic diffusion mechanics \\
      \bottomrule
    \end{tabularx}
    \begin{tablenotes}[flushleft]
      \footnotesize
    \item Target \& Scope: Seismotectonic and statistical criteria
      distinguishing elastostatic aftershock cascades from
      fluid-driven earthquake swarms.
    \item Notice: These are schematic literature interpretations, not
      deterministic causal classification rules. Branching ratio,
      magnitude distribution, and migration need separate estimation.
    \item Evidence Anchor: Grounded in rate-and-state friction
      mechanics and diffusion models. No regional catalogue benchmark
      was verified in this review.
    \end{tablenotes}
  \end{threeparttable}
\end{table}

% ==============================================================================
% SECTION 5: INFORMATION THEORY, THERMODYNAMICS, AND UNIFIED
% GENERATIVE FORMULATION
% ==============================================================================
\section{Information Theory, Thermodynamics, and Unified Generative
Formulation}
\label{sec:info_thermo_unified}

\subsection{Free Energy Landscape and Statistical Thermodynamics of Seismicity}
\label{subsec:free_energy_landscape}

We distinguish ideal-gas algebra from phenomenological physical analogies.
Neither specifies a physical earthquake temperature. Let $\mathcal{C} = \{(t_i,
m_i)\}_{i=1}^N \subset [0, T_{\mathrm{obs}}] \times [m_0, \infty)$
denote an earthquake catalogue recorded over observation horizon
$[0, T_{\mathrm{obs}}]$ with completeness magnitude $m_0 \in
\mathbb{R}$ (\textbf{Level~1: Raw Observations}).

\begin{definition}[Microscopic Effective Hamiltonian]
  \label{def:effective_hamiltonian}
  Let $\lambda_0 \equiv \qty{1}{event\per\day}$ be a dimensional
  reference scale. For any local rate or non-parametric density
  estimate $\hat{\lambda}(t) \in \mathbb{R}_{>0}$ (\textbf{Level~4:
  Model Estimate}), the \emph{effective Hamiltonian} at time $t \in
  [0, T_{\mathrm{obs}}]$ is the pointwise surprisal:
  \begin{equation}
    \mathcal{H}_{\mathrm{eff}}(t) \coloneqq
    -\ln\!\left(\frac{\hat{\lambda}(t)}{\lambda_0}\right).
    \label{eq:effective_hamiltonian}
  \end{equation}
  At discrete event times $t_i$, the microscopic event Hamiltonian
  is $\mathcal{H}_i \equiv \mathcal{H}_{\mathrm{eff}}(t_i) \coloneqq
  -\ln(\hat{\lambda}_i/\lambda_0)$.
\end{definition}

High-rate intervals ($\hat{\lambda} \gg \lambda_0$) generate deeply
negative potential wells ($\mathcal{H}_{\mathrm{eff}} \ll 0$),
defining energetically bound cluster cores (\textbf{Level~6:
Qualitative Interpretation}). Quiescent background regimes and
inter-cluster boundaries define potential barriers
($\mathcal{H}_{\mathrm{eff}} \gg 0$).

\begin{definition}[Uncertainty Proxy, Not a Physical Temperature]
  \label{prop:effective_temperature}
  Let $\kstar(t) \in \{3, \dots, k_{\max}\}$ denote the adaptive
  neighbourhood size selected by the $\pak$ estimator at time $t$.
  A fixed-rank delta-method benchmark motivates the dimensionless
  uncertainty proxy $T_{\mathrm{proxy}}(t)>0$:
  \begin{equation}
    T_{\mathrm{proxy}}(t) \coloneqq
    \frac{1}{\kstar(t) - 2}\,, \qquad \kstar(t) \ge \num{3}.
    \label{eq:effective_temperature}
  \end{equation}
  This is neither the exact fixed-rank variance $\psi_1(k)$ nor
  calibrated adaptive variance or thermal energy. Increasing $N$
  does not force $k^*\to\infty$ when ranks are capped.
\end{definition}

\begin{theorem}[Canonical and Grand Canonical Partition Functions]
  \label{thm:partition_functions}
  For deterministic Poisson intensity, conditional on a positive event
  count, unordered event locations have density $\lambda/\Lambda$,
  where $\Lambda =
  \int_0^{T_{\mathrm{obs}}} \lambda(s)\,\mathrm{d}s \in
  \mathbb{R}_{>0}$, the probability density of an event arriving at
  $t \in [0, T_{\mathrm{obs}}]$ matches the canonical
  Gibbs--Boltzmann distribution:
  \begin{equation}
    p(t)=\frac{\lambda(t)}{\Lambda}
    =\frac{\lambda_0e^{-\mathcal H_{\mathrm{eff}}(t)}}{\mathcal Z},
    \qquad \mathcal Z=\int_0^{T_{\mathrm{obs}}}
    e^{-\mathcal H_{\mathrm{eff}}(s)}\lambda_0\,\mathrm{d}s=\Lambda.
    \label{eq:canonical_boltzmann}
  \end{equation}
  In the grand canonical ensemble where total event count $n \in
  \mathbb{N}_0$ fluctuates, the grand partition function $\Xi$ for
  the Poisson process with chemical potential
  $\mu_{\mathrm{chem}}(t) \equiv \ln(\lambda(t)/\lambda_0) =
  -\mathcal{H}_{\mathrm{eff}}(t)$ is:
  \begin{equation}
    \Xi = \sum_{n=0}^{\infty} \frac{1}{n!} \left(
    \int_0^{T_{\mathrm{obs}}} \lambda(t)\,\mathrm{d}t \right)^n
    = \exp\!\left( \int_0^{T_{\mathrm{obs}}}
    \lambda(t)\,\mathrm{d}t \right) = e^{\Lambda}\,.
    \label{eq:grand_partition}
  \end{equation}
  The formal dimensionless free energy is $-\ln\mathcal Z=-\ln\Lambda$,
  with reference measure $\lambda_0\,\mathrm{d}t$. This location identity
  does not generally apply to history-dependent Hawkes intensities.
\end{theorem}

\begin{proposition}[Landau Field Theory of Swarm Onset]
  \label{prop:landau_field_theory}
  Following Landau's phenomenological theory of phase transitions
  \parencite{landau1937}, let $\mu_0 > 0$ denote the secular
  background rate. Define the scalar order parameter $\phi(t) \in
  [-\mu_0, \infty)$ by $\phi(t) \coloneqq \lambda(t) - \mu_0$. The free
  energy functional governing temporal cluster formation up to
  quartic order is:
  \begin{equation}
    \mathcal{F}[\phi] \coloneqq \int_0^{T_{\mathrm{obs}}} \left[
      \frac{\kappa}{2}\,\dot{\phi}(t)^2
      + \frac{a(t)}{2}\,\phi(t)^2 + \frac{c}{3}\,\phi(t)^3 +
    \frac{b}{4}\,\phi(t)^4 \right] \,\mathrm{d}t\,, \qquad
    \dot{\phi}(t) \equiv \frac{\mathrm{d}\phi}{\mathrm{d}t}\,,
    \label{eq:landau_functional}
  \end{equation}
  where $\kappa > 0$ is temporal gradient stiffness, $b > 0$ ensures
  stability, and $a(t) = a_0(t_c - t)$ ($a_0 > 0$) parameterises the
  distance from the critical onset time $t_c \in (0, T_{\mathrm{obs}})$.
  \begin{enumerate}[label=(\roman*),noitemsep]
    \item \textbf{Continuous (Second-Order) Swarm Onset:} For $c =
      0$, $a(t) > 0$ yields a single disordered ground state $\phi^*
      = 0$ (quiescent background). At $t > t_c$, $a(t) < 0$ induces
      spontaneous symmetry breaking with continuous order parameter growth:
      \begin{equation}
        \phi^*(t) \coloneqq \sqrt{\frac{-a(t)}{b}} = \sqrt{\frac{a_0(t -
        t_c)}{b}} \propto (t - t_c)^{1/2}, \qquad t \ge t_c.
        \label{eq:critical_exponent}
      \end{equation}
    \item \textbf{Discontinuous (First-Order) Mainshock Cascades:}
      For $c<0$, the unconstrained polynomial has coexistence at
      $a=2c^2/(9b)$ and $\phi=-2c/(3b)>0$. This is a phenomenological
      analogy, not a derived mainshock mechanism.
    \item \textbf{Mean-Field Limitation:} Neglecting gradients needs
      independent justification. Large $k^*$ is not a Ginzburg criterion.
      The negative branch for $c=0$ is admissible only while
      $\phi\ge-\mu_0$; the displayed positive branch does not
      demonstrate a seismic critical exponent.
  \end{enumerate}
\end{proposition}

\begin{table}[htbp]
  \centering
  \small
  \begin{threeparttable}
    \caption{Formal statistical-mechanics identities and proposed seismic
      analogies. The Poisson partition identity is distinct from the
      phenomenological phase labels and uncertainty proxy. No physical
    temperature, seismic phase transition, or empirical exponent is verified.}
    \label{tab:thermo_correspondence}
    \begin{tabularx}{\linewidth}{>{\raggedright\arraybackslash}X
      >{\raggedright\arraybackslash}X >{\raggedright\arraybackslash}X}
      \toprule
      \textbf{Seismological Observable} & \textbf{Thermodynamic
      Counterpart} & \textbf{Mathematical Definition} \\
      \midrule
      Background seismicity ($\lambda \approx \mu_0$) & Disordered
      (paramagnetic) phase & $\phi(t) = 0$ \\
      Swarm excursion ($\lambda \gg \mu_0$) & Ordered (ferromagnetic)
      phase & $\phi(t) = \sqrt{-a/b} > 0$ \\
      Mainshock-aftershock cascade & First-order phase transition &
      Discontinuous jump via $c\,\phi^3$ \\
      Aseismic / fluid-driven onset & Second-order continuous
      transition & Power-law onset $\phi \sim (t - t_c)^{1/2}$ \\
      Inter-cluster quiescence / saddle & Domain wall / phase boundary
      & Local minimum of $\lambda(t)$, saddle $\hat{\lambda}_s$ \\
      Finite-sample precision $\kstar$ & Dimensionless uncertainty
      proxy & $T_{\mathrm{proxy}} = 1/(\kstar - 2)$ \\
      Local surprisal $-\ln(\lambda/\lambda_0)$ & Effective
      Hamiltonian $\mathcal{H}_{\mathrm{eff}}$ &
      $-\ln(\lambda/\lambda_0)$ \\
      Total integrated intensity $\Lambda$ & Canonical partition
      function $\mathcal{Z}$ & $\mathcal{Z}=\Lambda$ \\
      \bottomrule
    \end{tabularx}
    \begin{tablenotes}[flushleft]
      \footnotesize
    \item [Model / Hypothesis] Phase labels, $\phi=\lambda-\mu_0$,
      polynomial coefficients $a,b,c$, and critical time $t_c$ are
      phenomenological. $\sqrt{-a/b}$ requires $a<0$, $b>0$, $c=0$.
    \item [Derived Identity] $\Lambda=\int\lambda\,\mathrm{d}t$ is
      dimensionless and $\lambda_0>0$ is a same-unit reference rate.
      The partition identity uses deterministic Poisson intensity;
      $T_{\mathrm{proxy}}$ is not thermal energy or calibrated
      adaptive variance.
    \end{tablenotes}
  \end{threeparttable}
\end{table}

\subsection{Information Geometry of Point Processes and Divergence Measures}
\label{subsec:information_geometry}

Let $\mathcal{M}_S$ denote the statistical manifold of point-process
conditional intensity models parameterised by vector
$\bm{\theta} \in \Theta \subset \mathbb{R}^d$.

\begin{theorem}[Point-Process Kullback--Leibler Divergence]
  \label{thm:pp_kl_divergence}
  Let $\mathbb{P}_A$ and $\mathbb{P}_B$ be probability measures on
  the path space of marked point configurations on $[0,
  T_{\mathrm{obs}}]$, induced by intensities $\lambda_A(t \mid
  \mathcal{H}_t)$ and $\lambda_B(t \mid \mathcal{H}_t)$
  respectively. Assume compatible initial histories, identical mark
  laws, absolute continuity, and finite log-integrals. Writing
  $\lambda_A=\lambda_A(t\mid\mathcal H_t)$ and similarly for $B$:
  \begin{equation}
    \mathrm D_{\mathrm{KL}}(\mathbb P_A\parallel\mathbb P_B)
    =\mathbb E_A\!\int_0^{T_{\mathrm{obs}}}
    \left[\lambda_A\ln\!\left(\frac{\lambda_A}{\lambda_B}\right)
    -\lambda_A+\lambda_B\right]\,\mathrm{d}t.
    \label{eq:pp_kl_divergence}
  \end{equation}
  For inhomogeneous Poisson processes with deterministic rates
  $\lambda_A(t), \lambda_B(t) \in \mathbb{R}_{>0}$, this evaluates to:
  \begin{equation}
    \mathrm{D}_{\mathrm{KL}}(\lambda_A \parallel \lambda_B) =
    \int_0^{T_{\mathrm{obs}}} \lambda_A(t)
    \ln\!\left(\frac{\lambda_A(t)}{\lambda_B(t)}\right)\,\mathrm{d}t -
    \int_0^{T_{\mathrm{obs}}} \bigl[\lambda_A(t) -
    \lambda_B(t)\bigr]\,\mathrm{d}t \;\ge\; 0.
    \label{eq:poisson_kl_explicit}
  \end{equation}
\end{theorem}

\begin{definition}[Jeffreys Divergence for Deterministic Rates]
  \label{def:jeffreys_divergence}
  The symmetrised Kullback--Leibler divergence across boundary
  window $\Delta t \in \mathbb{R}_{>0}$ is:
  \begin{equation}
    J(\lambda_A, \lambda_B) \coloneqq
    \mathrm{D}_{\mathrm{KL}}(\lambda_A \parallel \lambda_B) +
    \mathrm{D}_{\mathrm{KL}}(\lambda_B \parallel \lambda_A) =
    \int_{\Delta t} \bigl[\lambda_A(t) - \lambda_B(t)\bigr]
    \ln\!\left(\frac{\lambda_A(t)}{\lambda_B(t)}\right)\,\mathrm{d}t.
    \label{eq:jeffreys_divergence}
  \end{equation}
  For piecewise-constant rate segments $\lambda_A, \lambda_B \in
  \mathbb{R}_{>0}$, $J_{AB} = (\lambda_A - \lambda_B)(\ln \lambda_A
  - \ln \lambda_B)\,\Delta t \ge 0$. Statistically, $J_{AB}$ is
  a statistical discrepancy across a boundary window. Its nonnegativity
  does not establish physical entropy production or time-reversal breaking.
\end{definition}

\begin{theorem}[Flatness of the Fisher--Rao Metric in Log-Density Coordinates]
  \label{thm:fisher_rao_flatness}
  Within the differential-geometric framework of information geometry
  \parencite{amari1985}, let $\rho \in \mathbb{R}_{>0}$ denote local
  Poisson intensity. In the fixed-rank experiment
  $V_k\sim\mathrm{Gamma}(k,\text{rate }\rho)$, with the same known
  $k$ throughout, $g_{\rho\rho}=k/\rho^2$. Under
  $\phi=\ln(\rho/\rho_0)$ for a same-unit reference intensity
  $\rho_0>0$, the Fisher--Rao line
  element transforms to:
  \begin{equation}
    \mathrm{d}s_{\mathrm{FR}}^2 = g_{\rho\rho}\,\mathrm{d}\rho^2 =
    \frac{k}{\rho^2}\,(\rho\,\mathrm{d}\phi)^2 = k\,\mathrm{d}\phi^2.
    \label{eq:fisher_metric_flat}
  \end{equation}
  Consequently, the Christoffel symbols vanish identically
  ($\Gamma^\phi_{\phi\phi} = 0$), the Riemannian curvature tensor is
  zero ($R^\phi{}_{\phi\phi\phi} \equiv 0$), and Fisher--Rao
  geodesics connecting densities $\rho_A$ and $\rho_B$ are straight
  lines in log-density coordinates:
  \begin{equation}
    \phi(s) = (1 - s)\phi_A + s\,\phi_B, \quad s \in [0, 1], \qquad
    d_{\mathrm{FR}}(\rho_A, \rho_B) \coloneqq \sqrt{k}\,\bigl|\ln \rho_A -
    \ln \rho_B\bigr|.
    \label{eq:geodesic_distance}
  \end{equation}
  For fixed volume $V$, instead $K\sim\mathrm{Poisson}(\rho V)$,
  giving $g_{\rho\rho}=V/\rho$ and
  $d_{\mathrm{FR}}=2\sqrt V|\sqrt{\rho_A}-\sqrt{\rho_B}|$.
  \Cref{fig:supp_sampling_experiments} distinguishes these experiments.
\end{theorem}

\subsection{The Unified Generative Intensity Formulation}
\label{subsec:unified_model_spec}

To resolve the physical superposition of secular tectonics,
transient fluid driving, and stress-mediated aftershocks, we
propose the Unified Generative Model. Component labels are model
interpretations, not identified physical causes; this EM/model-selection
workflow is not verified executable behavior in the inspected module.

\begin{definition}[Unified Conditional Intensity]
  \label{def:unified_conditional_intensity}
  Let $\bm{\Theta} \in \bm{\Omega} \subset
  \mathbb{R}^{d}$ denote the complete parameter vector. The unified
  conditional intensity function for $t \in [0, T_{\mathrm{obs}}]$ is:
  \begin{equation}
    \lambda(t \mid \mathcal{H}_t; \bm{\Theta}) \;\equiv\;
    \underbrace{\mu_0(t)}_{\text{tectonic background}} \;+\;
    \underbrace{\sum_{k=1}^{K}
    A_k\,\phi_k(t;\,\bm{\theta}_k)}_{\text{exogenous swarm
    transients}} \;+\; \underbrace{\sum_{t_j < t}
    \kappa(m_j)\,\psi(t - t_j)}_{\text{endogenous ETAS triggering}}\,,
    \label{eq:unified_conditional_intensity}
  \end{equation}
  subject to parameter domains:
  \begin{align}
    \mu_0(t) &\equiv \sum_{j=0}^{J} \mu_j\,\mathbf{1}_{(\tau_j,
    \tau_{j+1}]}(t), \quad \mu_j \in \mathbb{R}_{>0},\; 0 = \tau_0 <
    \dots < \tau_{J+1} = T_{\mathrm{obs}}, \label{eq:background_spec} \\
    \phi_k(t;\,\bm{\theta}_k) &\equiv \frac{1}{(t -
    t_0^{(k)})\,\sigma_k\sqrt{2\pi}} \exp\!\left( -\frac{\bigl(\ln(t
      - t_0^{(k)}) - \nu_k\bigr)^2}{2\,\sigma_k^2}
    \right)\mathbf{1}_{(t_0^{(k)}, \infty)}(t), \label{eq:envelope_spec} \\
    \kappa(m) &\equiv K_0 \cdot 10^{\alpha(m - m_0)}, \quad K_0 \in
    \mathbb{R}_{>0},\; \alpha \in (0, b), \label{eq:etas_productivity} \\
    \psi(\tau) &\equiv \frac{(p - 1)\,c^{p-1}}{(\tau +
    c)^p}\,\mathbf{1}_{(0, \infty)}(\tau), \quad c \in
    \mathbb{R}_{>0},\; p \in (1, \infty), \label{eq:etas_kernel}
  \end{align}
  with amplitude $A_k \in \mathbb{R}_{>0}$, onset time $t_0^{(k)}
  \in [0, T_{\mathrm{obs}})$, log-location $\nu_k \in \mathbb{R}$,
  and log-scale $\sigma_k \in \mathbb{R}_{>0}$.
\end{definition}

\begin{proposition}[Analytical Point-Process Log-Likelihood and Compensator]
  \label{prop:analytical_loglik}
  For catalogue $\mathcal{C} = \{(t_i, m_i)\}_{i=1}^N$, the exact
  log-likelihood $\ln \mathcal{L}(\bm{\Theta}) =
  \sum_{i=1}^N \ln \lambda(t_i \mid \mathcal{H}_{t_i}) -
  \Lambda(T_{\mathrm{obs}})$ evaluates analytically via the
  compensator decomposition:
  \begin{equation}
    \begin{aligned}
      \Lambda(T_{\mathrm{obs}})
      &=\sum_{j=0}^{J}\mu_j(\tau_{j+1}-\tau_j) \\
      &\quad+\sum_{k=1}^{K}A_k\Phi_{\mathrm{LN}}\!
      \left(\ln(T_{\mathrm{obs}}-t_0^{(k)});\nu_k,\sigma_k\right) \\
      &\quad+\sum_{j=1}^{N}\kappa(m_j)
      \left[1-\left(\frac{c}{T_{\mathrm{obs}}-t_j+c}\right)^{p-1}\right],
    \end{aligned}
    \label{eq:analytical_compensator}
  \end{equation}
  where $\Phi_{\mathrm{LN}}(x; \nu, \sigma) = \frac{1}{2}\bigl[1 +
  \mathrm{erf}\bigl((x - \nu)/(\sigma\sqrt{2})\bigr)\bigr]$ is the
  Gaussian CDF evaluated at log-delay $x$, not the log-normal CDF of $x$.
  Assume empty pre-observation history and onsets in the stated interval;
  general onsets require both upper and lower CDF limits. Time logarithms
  use a fixed reference time unit, as in the supplementary derivation.
\end{proposition}

\subsection{Two-Stage Clustering Algorithm and Joint Model Selection}
\label{subsec:two_stage_clustering}

Parameter estimation and structural identification of the $K$
transient envelopes proceed via a two-stage sequential inference architecture.

\paragraph{Stage 1: Non-parametric Topological Segmentation ($\pak$).}
The raw catalogue times $\{t_i\}_{i=1}^N$ are evaluated by the
$\pak$ algorithm (\textbf{Level~3: Deterministic Computations}). At
each event $t_i$, the optimal neighbourhood $\kstar(t_i)$ is
selected via the rank statistic in \cref{def:kstar_stopping}.
Directed ascending
gradient flow on the density graph assigns events to basins of
attraction. Saddle points $\hat{\lambda}_s$ connecting adjacent
peaks $A$ and B are assessed via the log-density $\Zscore$-score:
\begin{equation}
  Z_{AB} \coloneqq \frac{\min(\ln \hat{\lambda}_A, \ln \hat{\lambda}_B) -
  \ln \hat{\lambda}_s}{\sqrt{\Var(\ln \hat{\lambda}_A) + \Var(\ln
  \hat{\lambda}_B) + \Var(\ln \hat{\lambda}_s)}}\,.
  \label{eq:zscore_topology}
\end{equation}
In this proposed workflow, boundaries with $Z_{AB}\le Z_{\mathrm{crit}}$
(default $Z_{\mathrm{crit}} \coloneqq \num{3.0}$) are merged. Surviving saddles
partition the catalogue into $K_{\mathrm{init}}$ temporal support
intervals $\{(a_k, b_k)\}_{k=1}^{K_{\mathrm{init}}}$. This is an
uncertainty-scaled heuristic, not a calibrated significance test or
the operational ADP merge rule below.

\paragraph{Stage 2: Parametric EM Refinement and Three-Way
Stochastic Declustering.}
Given initial windows $(a_k, b_k)$, log-normal parameters are
initialised from empirical moments, choosing
$\hat t_0^{(k)}=a_k<\min_{t_i\in C_k}t_i$ to exclude $\ln0$:
$\hat{\nu}_k = \frac{1}{|C_k|}\sum_{t_i \in C_k}\ln(t_i - a_k)$,
$\hat{\sigma}_k^2 = \frac{1}{|C_k|}\sum_{t_i \in C_k}(\ln(t_i - a_k)
- \hat{\nu}_k)^2$, and $\hat{A}_k = |C_k|$. The \ac{EM} algorithm
evaluates three-way posterior responsibilities (\textbf{Level~4:
Model Estimates}):
\begin{equation}
  w_i^{(\mathrm{bg})} \coloneqq \frac{\mu_0(t_i)}{\lambda(t_i \mid
  \mathcal{H}_{t_i})}, \quad
  w_i^{(\mathrm{env}, k)} \coloneqq \frac{A_k\,\phi_k(t_i)}{\lambda(t_i \mid
  \mathcal{H}_{t_i})}, \quad
  w_i^{(\mathrm{trig})} \coloneqq \frac{\sum_{t_j <
  t_i}\kappa(m_j)\,\psi(t_i - t_j)}{\lambda(t_i \mid \mathcal{H}_{t_i})},
  \label{eq:three_way_responsibilities}
\end{equation}
satisfying $w_i^{(\mathrm{bg})} + \sum_{k=1}^K w_i^{(\mathrm{env},
k)} + w_i^{(\mathrm{trig})} \equiv 1$, updating parameters via
weighted maximum likelihood.

\begin{definition}[Joint $\mdl$--$\pak$ Model Selection Criterion]
  \label{def:joint_mdl_pak}
  Under Rissanen's Minimum Description Length principle
  \parencite{rissanen1978}, for candidate model order $K \in \{0,
  \dots, K_{\mathrm{init}}\}$, let $d_K = 4K + 4 + (J + 1)$ denote the
  total free parameter count. The joint objective functional
  $\mathcal{J}(K)$ is:
  \begin{equation}
    \mathcal{J}(K) \coloneqq \underbrace{-\ln
      \mathcal{L}(\hat{\bm{\Theta}}_K) + \frac{d_K}{2}\ln
    N}_{\mathrm{MDL}(K) \text{ (Two-Part Code Length)}} \;+\;
    \underbrace{\gamma \sum_{s=1}^{\max(K-1,0)} \max\bigl(0,\;
    Z_{\mathrm{crit}} - Z_s\bigr)}_{\text{Topological Density Penalty}}\,,
    \label{eq:joint_mdl_pak}
  \end{equation}
  with coupling parameter $\gamma \coloneqq \frac{1}{2}\ln N$. The optimal
  envelope order is:
  \begin{equation}
    K^*=\argmin_{0\le K\le K_{\mathrm{init}}}\mathcal J(K).
  \end{equation}
  The parameter count assumes fixed background knots; free knots add
  parameters. At $K=0$, there is no envelope or topological split penalty.
\end{definition}

\begin{remark}[Model-Selection Consistency Remains Unproved]
  \label{thm:joint_consistency}
  A desired target for true envelope count $K_{\mathrm{true}}$ is:
  \begin{equation}
    \mathbb P(K^*=K_{\mathrm{true}})\longrightarrow1.
    \label{eq:consistency_result}
  \end{equation}
  This is consistency in probability, not strong consistency. Log-normal
  envelopes lack compact support; vanishing widths can make likelihood
  unbounded, while extra components cause boundary and identifiability
  problems. A proof needs constrained widths, bounded parameter domains,
  an explicit asymptotic experiment, identifiability, and selection-aware
  calibration. Classical Wilks arguments do not establish this target.
\end{remark}
\begin{landscape}
  \begin{table}[htbp]
    \centering
    \small
    \begin{threeparttable}
      \caption{Comparison of density methods, ETAS literature models, and
        the proposed unified temporal model. Additive intensity components
        allow model-based attribution, not unique physical causation.
        Density operations were source-inspected; no unified-model
      implementation, calibration, or benchmark run was verified.}
      \label{tab:model_comparison_chapter8}
      \begin{tabularx}{\linewidth}{>{\raggedright\arraybackslash}X
          >{\raggedright\arraybackslash}X >{\raggedright\arraybackslash}X
        >{\raggedright\arraybackslash}X}
        \toprule
        \textbf{Attribute} & \textbf{Poisson / $\pak$} & \textbf{Pure
        $\etas$} & \textbf{Unified Generative Model} \\
        \midrule
        Intensity & $\hat\lambda(t)$ & $\mu_0+\lambda_{\mathrm{trig}}$ &
        $\mu_0(t)+\lambda_{\mathrm{env}}+\lambda_{\mathrm{trig}}$ \\
        Epistemic Basis & Level 3/4 (Non-parametric) & Level 4/5
        (Mechanistic branching) & Level 4/5 (Joint superposition) \\
        Driving Mechanism & Non-specific rate peaks & Endogenous stress
        transfer interpretation & Envelope forcing $+$ triggering;
        forcing is not uniquely identified as fluid \\
        Declustering Mode & None (Binary hard threshold) & Two-way
        ($w_{\mathrm{bg}}, w_{\mathrm{trig}}$) & Three-way
        ($w_{\mathrm{bg}}, w_{\mathrm{env}}, w_{\mathrm{trig}}$) \\
        Boundary Physics & Density saddle $\Zscore$-score & Offspring
        cluster graph & Jeffreys statistical discrepancy $J_{AB}$ \\
        Model Selection & Non-parametric persistence & Likelihood ratio
        / $\mathrm{AIC}$ & Joint $\mathrm{MDL}$--$\pak$ criterion
        $\mathcal{J}(K)$ \\
        Information Benchmark & Fixed-rank Fisher information &
        Regularity-dependent information bounds & Not established
        for the proposed mixture-selection problem \\
        \bottomrule
      \end{tabularx}
      \begin{tablenotes}[flushleft]
        \footnotesize
      \item [Literature Baseline / Proposal] ETAS is epidemic-type
        aftershock sequencing; MDL is minimum description length;
        AIC is Akaike information criterion. $\lambda_{\mathrm{env}}
        =\sum_kA_k\phi_k$ and $\lambda_{\mathrm{trig}}
        =\sum_{t_j<t}\kappa(m_j)\psi(t-t_j)$ are rates, while
        $w$ denotes conditional component responsibilities.
      \end{tablenotes}
    \end{threeparttable}
  \end{table}
\end{landscape}
% ==============================================================================
% SECTION 6: ADVANCED DENSITY PEAKS AND VECTORIZED OPTIMIZATIONS (ADP++)
% ==============================================================================
\section{Advanced Density Peaks and Vectorized Algorithmic
Optimizations (ADP++)}
\label{sec:adppp_optimizations}

The non-parametric identification of cluster cores, borders, and
halos on arbitrary manifolds builds upon the Advanced Density Peaks
(\textsc{ADP} / \textsc{DPA}) algorithm
\parencite{rodriguez2014,derrico2021}. However, standard scalar
implementations suffer from severe execution bottlenecks
($\mathcal{O}(N \bar{k})$ Python loops and redundant heap
traversals), restricting scalability on dense seismic catalogs ($N
\ge \num{1e5}$). This section presents the vectorized \textsc{ADP++}
algorithmic suite, proving convergence on its parent forest and describing
source-verified vectorized operations. Universal scalar parity and
numerical speedups require separate tests and timing records.

\subsection{Standard Advanced Density Peaks Pipeline}
\label{subsec:adp_standard}

Given dataset $\mathcal{D} = \{\mathbf{x}_i\}_{i=1}^N \subset
\mathbb{R}^D$ with \pak{} log-densities $\hat{u}_i \coloneqq \ln
\hat{\rho}_i$ and standard errors $\hat{\sigma}_i$ (\textbf{Level~4:
Model Estimates}):
\begin{enumerate}[noitemsep,topsep=2pt]
  \item \textbf{Peak Identification:} Candidate modes are identified
    using the default score $g_i=\hat u_i-\hat\sigma_i$. Point
    $\mathbf{x}_i$ is a
    candidate peak if $g_i \ge g_j$ for all $j \in \mathcal{N}_i^{\kstar(i)}$.
  \item \textbf{Basin Assignment:} Non-peak points are assigned to
    cluster basins by directed steepest ascent along the density
    gradient graph $\mathcal{F}$.
  \item \textbf{Saddle-Point Detection:} For adjacent clusters $C_c,
    C_{c'}$, border points are pairs $(i, j)$ such that $i \in C_c$,
    $j \in C_{c'}$, and $j \in \mathcal{N}_i^{\kstar(i)}$. The
    ADP candidates use each point's first cross-cluster neighbour,
    defining edge set $\mathcal E_{cc'}^{(\mathrm{first})}$. Its
    saddle proxy is:
    \begin{equation}
      \hat u_{cc'}^\dagger=
      \max_{(i,j)\in\mathcal E_{cc'}^{(\mathrm{first})}}
      \frac{\hat u_i+\hat u_j}{2}.
      \label{eq:saddle_density_def}
    \end{equation}
  \item \textbf{Multimodality Merging:} Clusters $C_c$ and $C_{c'}$
    are merged if the density drop between peak and saddle is
    below the operational uncertainty margin for either peak:
    \begin{equation}
      \hat u_c-\hat u_{cc'}^\dagger
      <Z_{\mathrm{conf}}(\hat\sigma_c+\hat\sigma_{cc'}^\dagger)
      \quad\text{or the analogous condition for }c'.
      \label{eq:merge_condition}
    \end{equation}
    The selected edge has $\hat\sigma_{cc'}^\dagger
    =(\hat\sigma_i+\hat\sigma_j)/2$; default $Z_{\mathrm{conf}}=1.65$
    is an algorithm parameter, not a calibrated significance level.
    The retained centre is selected by normalized margins, not
    necessarily the elder peak of persistent homology.
  \item \textbf{Halo Filtering:} Points in $C_c$ with density below
    the cluster border threshold ($\hat{u}_i < \max_{c'}
    \hat{u}_{cc'}^\dagger$) are labeled as noise/halo
    ($\ell_i = -1$).
\end{enumerate}

\subsection{Vectorized ADP++ Algorithmic Architecture}
\label{subsec:adppp_vectorized}

The \textsc{ADP++} implementation eliminates interpreter overhead
through four fused vectorized transformations (\textbf{Level~3:
Deterministic Computations}):

\paragraph{Optimization 1: Fused Vectorized Union-Find (Steps 5 \& 6).}
Rather than traversing scalar path pointers, \textsc{ADP++}
constructs a vectorized parent link array $\mathbf{p} \in
\mathbb{Z}^N$ by broadcasting neighbor scores:
\begin{equation}
  p_i \coloneqq \argmax_{j \in \mathcal{N}_i^{\kstar(i)},\; g_j > g_i} g_j,
  \qquad (\text{with } p_i \coloneqq i \text{ if candidate peak}).
  \label{eq:vectorized_parent}
\end{equation}
Cluster basin assignment is executed via parallel pointer jumping:
\begin{equation}
  \mathbf{r}^{(t+1)} \coloneqq \mathbf{r}^{(t)}\bigl[\mathbf{r}^{(t)}\bigr],
  \qquad \mathbf{r}^{(0)} \coloneqq \mathbf{p}.
  \label{eq:pointer_jumping}
\end{equation}

\begin{theorem}[Equivalence and Convergence of Fused Union-Find]
  \label{thm:union_find_equivalence}
  The directed graph $\mathcal{G}_{\mathrm{ascent}} = (\mathcal{V},
  \mathcal{E})$ defined by directed edges $(i, p_i)$ is a collection
  of disjoint directed trees rooted at the candidate peaks. Pointer
  jumping in \cref{eq:pointer_jumping} converges to sequential traversal
  of the same parent forest in at most
  $\lceil\log_2\max(1,L_{\max})\rceil\le\lceil\log_2N\rceil$ parallel vector
  lookups, where $L_{\max}$ is the length of the longest gradient path.
\end{theorem}

\begin{proof}
  By definition, $p_i \ne i \implies g_{p_i} > g_i$. Hence
  $\mathcal{G}_{\mathrm{ascent}}$ is strictly acyclic. Each non-peak
  node has out-degree 1, and each peak has out-degree 0 after
  omitting self-loops.
  Thus $\mathcal{G}_{\mathrm{ascent}}$ is a directed forest rooted
  at peaks. After $t$ iterations of $\mathbf{r}^{(t+1)} =
  \mathbf{r}^{(t)}[\mathbf{r}^{(t)}]$, $\mathbf{r}^{(t)}[i]$ points
  to the ancestor at distance $\min(2^t,L(i))$ along the unique path from $i$
  to its root. When $2^t \ge L_{\max}$, $\mathbf{r}^{(t)}[i]$ is
  identical to the peak root, matching sequential path compression identically
  \parencite{tarjan1975}.
\end{proof}

Each vector update costs $\mathcal O(N)$ work; an already-rooted forest
needs zero updates. Scalar ADP additionally prunes centres and handles
ties through a different assignment path, so this theorem does not prove
general scalar/vectorized parity. \Cref{fig:supp_pointer_jumping} separates
ancestor doubling from the total-work bound.

\paragraph{Optimization 2: Vectorized Saddle-Point Detection (Step 7).}
Let $\mathbf{C}_{\mathrm{nb}} \in \mathbb{Z}^{N \times k_{\max}}$ be
the cluster label tensor of nearest neighbours. The boolean border
mask $\mathbf{B} \in \{0, 1\}^{N \times k_{\max}}$ is constructed
via broadcasting:
\begin{equation}
  B_{ij} = \mathbf{1}_{\{C_{\mathrm{nb}, ij} \ne \ell_i\}} \cdot
  \mathbf{1}_{\{j \le \kstar(i)\}}.
  \label{eq:border_tensor}
\end{equation}
The sequential first-hit condition (`break` upon encountering a
distinct cluster) is vectorized via:
\begin{equation}
  j_{\mathrm{first}}(i) = \argmax_{j \in \{1, \dots, k_{\max}\}}
  \bigl( B_{ij} \bigr).
  \label{eq:vectorized_first_hit}
\end{equation}
Points with $\max_j B_{ij} = 1$ are candidate border points. Sorting
candidate boundary edges in descending order of
$(\hat u_i+\hat u_{j_{\mathrm{first}}})/2$ and updating the symmetric saddle
matrix $\hat{\mathbf{U}}^\dagger \in \mathbb{R}^{C \times C}$ in
vectorized COO format executes in $\mathcal{O}(N k_{\max} + E\log
E)$ time, where $E \le N$ is the number of border points.

\paragraph{Optimization 3: Vectorized $k$-Peaks Density (Step 4d).}
Let $\mathbf{K}_{\mathrm{nb}} \in \mathbb{Z}^{N \times k_{\max}}$ be
the matrix of neighbor $\kstar$ values. The masked variance matrix
is evaluated via row-wise tensor reductions:
\begin{equation}
  \mathbf{\Delta}^2 = \bigl( \mathbf{K}_{\mathrm{nb}} -
  \mathbf{k}^* \mathbf{1}^\top \bigr)^2 \odot
  \mathbf{M}_{\kstar}, \qquad \bar{\sigma}_i^2 =
  \frac{1}{\kstar(i)}\sum_{j=1}^{k_{\max}} \Delta_{ij}^2,
  \label{eq:vectorized_kpeaks}
\end{equation}
The helper uses raw $k^*$ values as proxies, not log-density estimates,
and still allocates arrays. The public density dispatcher currently aliases
\texttt{kpeaks} to \texttt{kstarNN}, rather than invoking this helper.

\paragraph{Optimization 4: Vectorized Multimodality Test (Step 8).}
Let $\hat{\mathbf{u}}_{\mathrm{peak}} \in \mathbb{R}^C$ and
$\hat{\bm{\sigma}}_{\mathrm{peak}} \in \mathbb{R}^C$ be
cluster peak density and error vectors. Outer-product broadcasting
constructs the candidate margin and uncertainty tensors:
\begin{equation}
  \mathbf{A}^{(1)} = \hat{\mathbf{u}}_{\mathrm{peak}}
  \mathbf{1}^\top - \hat{\mathbf{U}}^\dagger, \quad
  \mathbf{A}^{(2)} = \mathbf{1}
  \hat{\mathbf{u}}_{\mathrm{peak}}^\top - \hat{\mathbf{U}}^\dagger, \quad
  \mathbf{E}^{(1)} =
  Z_{\mathrm{conf}}\bigl(\hat{\bm{\sigma}}_{\mathrm{peak}}
  \mathbf{1}^\top + \hat{\mathbf{\Sigma}}^\dagger\bigr), \quad
  \mathbf{E}^{(2)} =
  Z_{\mathrm{conf}}\bigl(\mathbf{1}\hat{\bm{\sigma}}_{\mathrm{peak}}^\top
  + \hat{\mathbf{\Sigma}}^\dagger\bigr).
  \label{eq:outer_product_merging}
\end{equation}
The boolean merge eligibility tensor is:
\begin{equation}
  \mathbf{M}_{\mathrm{elig}} = \bigl(\mathbf{A}^{(1)} <
  \mathbf{E}^{(1)}\bigr) \lor \bigl(\mathbf{A}^{(2)} < \mathbf{E}^{(2)}\bigr).
  \label{eq:merge_eligibility}
\end{equation}
Selecting the highest finite positive saddle value over eligible
upper-triangular pairs, with no merge when the eligible set is empty, and
updating border densities via in-place row-wise maximization
$\hat{\mathbf{U}}^\dagger[c, :] \leftarrow
\max(\hat{\mathbf{U}}^\dagger[c, :], \hat{\mathbf{U}}^\dagger[c',
:])$ avoids re-evaluating the $N$-point point cloud.
\begin{landscape}
  \begin{table}[htbp]
    \centering
    \small
    \caption{\textbf{Theoretical Computational Complexity: Standard
      ADP vs.\ Vectorized ADP++}. Target/Scope: Algorithmic asymptotic
      work bounds and unverified draft timings. Vectorization can remove
      Python-loop overhead but does not prove bitwise parity or a speedup.
      Bounds exclude neighbour construction; the draft timing setup
    $N=\num{1e5}$, $k_{\max}=80$, $C=50$ has no inspected run record.}
    \label{tab:complexity_comparison}
    \begin{threeparttable}
      \begin{tabularx}{\linewidth}{>{\raggedright\arraybackslash}X
          >{\raggedright\arraybackslash}X >{\raggedright\arraybackslash}X
        >{\raggedright\arraybackslash}X}
        \toprule
        \textbf{Pipeline Stage} & \textbf{Standard \textsc{ADP}} &
        \textbf{Vectorized \textsc{ADP++}} & \textbf{Unverified
        Draft Speedup} \\
        \midrule
        Step 5 \& 6: Ascent \& Basin Assignment & $\mathcal{O}(N
        \bar{k})$ scalar loops & $\mathcal{O}(N\log N+NC)$ including
        basin scans
        & $\num{140}\times$ \\
        Step 7: Saddle Detection & $\mathcal{O}(N k_{\max})$ scalar
        checks & $\mathcal{O}(N k_{\max} + E\log E)$ tensor ops &
        $\num{85}\times$ \\
        Step 4d: $k$-Peaks Variance & $\mathcal{O}(N k_{\max})$ loop
        allocations & $\mathcal{O}(N k_{\max})$ masked SIMD &
        $\num{45}\times$ \\
        Step 8: Multimodality Test & $\mathcal{O}(C^2 \times
        N_{\mathrm{iter}})$ re-scans & $\mathcal{O}(C^2 +
        C_{\mathrm{merges}} C^2)$ pair scans & $\num{310}\times$ \\
        \textbf{Total End-to-End Pipeline} & $\mathcal{O}(N \bar{k} +
        C^2 N_{\mathrm{iter}})$ & $\mathcal{O}(N\log N+Nk_{\max})$
        plus $\mathcal{O}(C^3+NC)$ worst-case work &
        $\mathbf{\numrange{120}{450}\times}$ \\
        \bottomrule
      \end{tabularx}
      \begin{tablenotes}[flushleft]
        \footnotesize
      \item [Derived Work Bound] $N$ is the point count, $C$ the
        initial basin count, $E\le N$ the first-hit edge count,
        $\bar k$ the mean selected rank, and $k_{\max}$ the rank cap.
        $N_{\mathrm{iter}}$ and $C_{\mathrm{merges}}$ count merge iterations.
      \item [Unverified Draft] Speedup values are retained for review,
        not reported as measurements. The $k$-peaks row describes
        an optional helper, not the current default dispatch.
      \end{tablenotes}
    \end{threeparttable}
  \end{table}
\end{landscape}
% ==============================================================================
% SECTION 7: CLUSTER VALIDATION VIA DENSITY SEPARATION SCORE AND
% TOPOLOGICAL PERSISTENCE
% ==============================================================================
\section{Cluster Validation via Density Separation Score and
Topological Persistence}
\label{sec:ds_score_persistence}

\subsection{The PAk Density Separation Score (DS-Score)}
\label{subsec:ds_score_formulation}

To validate clusters on complex manifolds without Euclidean
convexity assumptions, we formalize the \pak{} Density Separation
Score (DS-Score) \parencite{derrico2021}.

\begin{definition}[Pairwise Cluster Density Separation Score]
  \label{def:pairwise_ds_score}
  Let $C_c$ and $C_{c'}$ be adjacent clusters with estimated peak
  log-densities $\hat{u}_c, \hat{u}_{c'}$, peak uncertainties
  $\hat{\sigma}_c, \hat{\sigma}_{c'}$, inter-cluster saddle
  log-density $\hat{u}_{cc'}^\dagger$, and saddle uncertainty
  $\hat{\sigma}_{cc'}^\dagger$. The pairwise Density Separation
  Score $Z_{cc'}$ is defined as:
  \begin{equation}
    Z_{cc'} \coloneqq \frac{\min(\hat{u}_c, \hat{u}_{c'}) -
    \hat{u}_{cc'}^\dagger}{\sqrt{\hat{\sigma}_c^2 +
    \hat{\sigma}_{c'}^2 + (\hat{\sigma}_{cc'}^\dagger)^2}}\,.
    \label{eq:pairwise_zscore}
  \end{equation}
  This is the operational denominator in the inspected code.
  A lower-peak-only alternative under a prespecified lower peak and
  independence approximation is:
  \begin{equation}
    Z_{cc'}^{(\mathrm{lower})} \coloneqq \frac{\min(\hat{u}_c, \hat{u}_{c'}) -
    \hat{u}_{cc'}^\dagger}{\sqrt{\hat{\sigma}_{\min(c, c')}^2 +
    (\hat{\sigma}_{cc'}^\dagger)^2}}\,.
    \label{eq:pairwise_zscore_robust}
  \end{equation}
\end{definition}

\begin{definition}[Global Aggregated PAk-DS Score]
  \label{def:global_ds_score}
  Let $\mathcal{A} \subset \{1, \dots, K\} \times \{1, \dots, K\}$
  denote the adjacency graph of clusters sharing border points. The
  global score $S_{\mathrm{PAk}}$ is the population-weighted
  arithmetic mean across all adjacent unordered pairs:
  \begin{equation}
    S_{\mathrm{PAk}} \coloneqq \frac{\sum_{(c, c') \in \mathcal{A}}
    w_{cc'} \, Z_{cc'}}{\sum_{(c, c') \in \mathcal{A}} w_{cc'}},
    \qquad w_{cc'} \coloneqq \sqrt{|C_c| \cdot |C_{c'}|}\,.
    \label{eq:global_pak_ds_score}
  \end{equation}
  If no saddle points exist (completely disconnected clusters in
  metric space), the score evaluates via virtual saddle fallback:
  \begin{equation}
    \hat{u}_{cc'}^\dagger \coloneqq \min_{\mathbf{x}_i \in \mathcal{D}}
    \hat{u}_i, \qquad \hat{\sigma}_{cc'}^\dagger=\hat\sigma_{i_{\min}},
    \quad i_{\min}=\argmin_i\hat u_i.
    \label{eq:virtual_saddle}
  \end{equation}
  PAk-DS uses all cross-cluster edges within $k_{\max}$ for adjacency,
  not ADP's first-hit edges. It maximizes edge minima
  $\min(\hat u_i,\hat u_j)$ and uses the selected edge error
  $\sqrt{\hat\sigma_i^2+\hat\sigma_j^2}$. A single cluster returns
  \texttt{None}; a virtual saddle is a fallback, not an observed saddle.
\end{definition}

\subsection{Topological Persistence Calibration and Morse Theory}
\label{subsec:topological_persistence}

The mathematical foundation of the DS-Score is rooted in Morse
theory and persistent homology (0-homology $\mathcal{H}_0$)
\parencite{edelsbrunner2002,chazal2013}.

\begin{theorem}[0-Homology Superlevel Filtration and Elder Rule]
  \label{thm:persistence_elder_rule}
  Let $f:\mathcal M\to\mathbb R$ be a positive Morse density on a
  compact boundary-free $d$-manifold, with distinct critical values.
  Consider the superlevel filtration $L(\lambda) \coloneqq
  \{\mathbf{x} \in \mathcal{M} : f(\mathbf{x}) \ge \lambda\}$ as
  threshold $\lambda$ decreases from $+\infty$ to $-\infty$:
  \begin{enumerate}[noitemsep,topsep=2pt]
    \item \textbf{Birth:} A new connected component of $L(\lambda)$
      is born at Morse index-$d$ critical points (local maxima /
      peaks) at birth value $b_c = f(\mathbf{x}_c^{\mathrm{peak}})$.
    \item \textbf{Death:} As $\lambda$ drops to index-$(d-1)$ critical
      points (saddle points $\mathbf{x}_{cc'}^\dagger$), two
      disjoint connected components $C_c$ and $C_{c'}$ merge at
      death value $d_{cc'} = f(\mathbf{x}_{cc'}^\dagger)$.
    \item \textbf{Elder Rule:} By the elder rule of persistent
      homology, the component with the higher birth value
      ($b_{\mathrm{elder}} = \max(b_c, b_{c'})$) survives, while the
      younger component ($b_{\mathrm{younger}} = \min(b_c, b_{c'})$) dies.
  \end{enumerate}
  The raw topological persistence of the younger component is
  $\pi_{cc'} \coloneqq \min(b_c, b_{c'}) - d_{cc'}$.
  These are conventional indices of $f$ on a $d$-manifold; indices
  0 and 1 instead refer to $-f$. A log transform preserves pairings,
  not lifetime values, as shown in \cref{fig:supp_persistence_pairing}.
\end{theorem}

\begin{remark}[Uncertainty-Scaled Persistence and Calibration Limits]
  \label{cor:calibrated_persistence}
  The pairwise \pak{} Density Separation Score $Z_{cc'}$ is the
  uncertainty-scaled log-density persistence proxy:
  \begin{equation}
    Z_{cc'}^{(\mathrm{lower})}=\frac{\hat\pi_u}{\sigma_{\mathrm{total}}},
    \qquad
    \sigma_{\mathrm{total}} \coloneqq \sqrt{\hat{\sigma}_{\min}^2 +
    (\hat{\sigma}^\dagger)^2}.
    \label{eq:calibrated_persistence}
  \end{equation}
  Here $\hat\pi_u=\min(\hat u_c,\hat u_{c'})-\hat u_{cc'}^\dagger$.
  Only for a prespecified contrast with a justified central limit
  theorem and consistent covariance-aware standard error would a
  boundary-null statistic $Z$ satisfy:
  \begin{equation}
    Z\xrightarrow{\mathcal D}\mathcal N(0,1),\qquad
    p_{\mathrm{nom}}=1-\Phi(Z).
    \label{eq:zscore_pvalue}
  \end{equation}
  At $Z=3$, the nominal one-sided tail is $\num{1.3499e-3}$ and the
  two-sided tail is $\num{2.6998e-3}$. Neither is a verified PAk-DS
  $p$-value. Selected extrema, adaptive ranks, and shared neighbours need
  calibration; false-discovery control also needs a multiplicity procedure.
\end{remark}

\begin{proposition}[Thermodynamic Kramers Escape Time Interpretation]
  \label{prop:kramers_escape}
  In statistical thermodynamics, the escape rate of a Brownian
  particle in an independently specified overdamped diffusion with
  energy potential $U$, noise energy $\epsilon=k_{\mathrm B}T>0$,
  and a large nondegenerate barrier is governed by the
  Arrhenius--Kramers rate \parencite{kramers1940}:
  \begin{equation}
    \tau_{\mathrm{escape}}\sim\tau_0e^{\Delta U/\epsilon}
    =\tau_0\left(
    \frac{\hat{\rho}_c}{\hat{\rho}_{cc'}^\dagger} \right)^{\beta}.
    \label{eq:kramers_time}
  \end{equation}
  The density-ratio equality additionally postulates
  $U=-E_0\ln(\rho/\rho_0)$ and $\beta=E_0/\epsilon$, with $E_0>0$.
  Mobility and curvature determine the prefactor. Statistical errors
  are not thermal noise; high PAk-DS alone implies neither this
  dynamics nor a physical residence time.
\end{proposition}

\subsection{Theoretical Comparison with Classical Metrics}
\label{subsec:cvi_theoretical_comparison}

\Cref{tab:cvi_comparison} systematically compares the
mathematical and structural properties of the \pak{}-DS Score
against classical validity indices.
\begin{landscape}
  \begin{table}[htbp]
    \centering
    \small
    \begin{threeparttable}
      \caption{\textbf{Theoretical Comparison: PAk-DS Score
        vs.\ Classical Cluster Validity Indices}. Target/Scope:
        Structural and geometric comparison of validity metrics on
        non-convex supports. Density contrasts address a different
        question from distance or centroid criteria, but adaptive
        geometry and selected extrema remain limitations. This is a
      mathematical comparison, not an empirical superiority result.}
      \label{tab:cvi_comparison}
      \begin{tabularx}{\linewidth}{l >{\raggedright\arraybackslash}X
        >{\raggedright\arraybackslash}X >{\raggedright\arraybackslash}X}
        \toprule
        \textbf{Metric / Index} & \textbf{Geometric Basis} &
        \textbf{Manifold Topology} & \textbf{Failure Mode on Seismicity} \\
        \midrule
        \textbf{Silhouette \parencite{rousseeuw1987}} &
        Point-to-cluster mean distance & Depends on chosen distance &
        May mismatch density connectivity; full pair distances cost
        $\mathcal O(N^2)$ \\
        \textbf{Calinski--Harabasz \parencite{calinski1974}} & Scatter
        matrix trace ratio & Centroid-based convex hulls & Centroids
        can lie off-manifold; sensitive to scatter and outliers \\
        \textbf{Davies--Bouldin \parencite{davies1979}} & Centroid
        distance vs.\ dispersion & Hyperspherical clusters & Fails on
        interleaved or branching fault networks \\
        \textbf{Dunn Index \parencite{dunn1974}} & Min inter-cluster /
        max diameter & Extreme pair distances & Pathologically
        sensitive to single outlier events \\
        \textbf{PAk-DS Score (Ours)} & \textbf{Morse saddle density
        contrast} & \textbf{Arbitrary non-convex $\mathcal{M}$} &
        \textbf{Selection, boundary bias, and uncalibrated extrema} \\
        \bottomrule
      \end{tabularx}
      \begin{tablenotes}[flushleft]
        \footnotesize
      \item Target \& Scope: Structural, mathematical, and geometric
        comparison of validity metrics on arbitrary $d$-dimensional
        Riemannian submanifolds.
      \item Notice: [Analyst Interpretation] Distance and centroid
        indices do not enforce convexity, but may disagree with density
        connectivity. PAk-DS uses selected neighbour-edge proxies,
        not exact continuous Morse saddles.
      \item Evidence Anchor: [Derived Definition] Displayed formulas
        and inspected PAk-DS code; no comparative benchmark result verified.
      \end{tablenotes}
    \end{threeparttable}
  \end{table}
\end{landscape}
% ==============================================================================
% SECTION 8: EMPIRICAL VALIDATION ON SYNTHETIC BENCHMARKS AND
% REGIONAL SEISMICITY
% ==============================================================================
\section{Synthetic Benchmark Drafts and Validation Requirements}
\label{sec:clustering_benchmarks}

\begin{remark}[Unverified Numerical Draft]
  The scores, timings, uncertainty intervals, ARIs, and Monte Carlo
  outcomes retained below have no inspected run records, seeds, or
  versioned results. They are unverified draft claims, not demonstrated
  results. Normal tails are not calibrated score $p$-values. Inspected
  tests provide examples, not provenance for these numbers.
\end{remark}

\subsection{Canonical Two-Dimensional Benchmarks}
\label{subsec:canonical_benchmarks}

To verify mathematical sensitivity, the \pak{}-DS Score and
\textsc{ADP++} pipeline were evaluated on five canonical synthetic
test topologies in the draft ($N=\num{2000}$ each), with numerical
outcomes pending reproducible validation:
\begin{enumerate}[label=(\roman*),noitemsep]
  \item \textbf{Well-Separated Gaussians:} Two bivariate isotropic
    Gaussians ($\mu_1 = \num{-3}, \mu_2 = \num{+3}, \sigma =
    \num{1.0}$). Evaluated
    score: $S_{\mathrm{PAk}} = \num{4.12} \pm \num{0.18}$ ($p <
    \num{1e-4}$). Clear two-cluster topology.
  \item \textbf{Overlapping Gaussians:} Centers shifted to $\mu_1 =
    \num{-1}, \mu_2 = \num{+1}$. Evaluated score: $S_{\mathrm{PAk}}
    = \num{0.59}
    \pm \num{0.08}$ ($p \approx \num{0.27}$). Unsupervised merging
    correctly collapses the distribution into a single cluster ($K = \num{1}$).
  \item \textbf{Uniform Poisson Noise:} Points sampled uniformly on
    $[0, 1]^2$. Evaluated score: $S_{\mathrm{PAk}} = \num{0.01} \pm
    \num{0.02}$ ($p \approx \num{0.49}$). Dynamic range between
    well-separated and noise: $\num{4.12} / \num{0.01} = \num{412}\times$.
  \item \textbf{Single Gaussian Mode:} Unimodal baseline. Evaluates
    to \texttt{None} in the inspected code (no inter-cluster
    saddles), triggering the single-cluster indicator.
  \item \textbf{Noise-Injected Separated Modes:} Two Gaussians
    corrupted by \qty{20}{\percent} uniform background noise. Halo filtering
    correctly assigns background points to $\ell_i = -1$; core
    separation score remains invariant at $S_{\mathrm{PAk}} =
    \num{4.13} \pm \num{0.21}$.
\end{enumerate}

\subsection{Comprehensive 21-Manifold Benchmark Suite}
\label{subsec:manifold_benchmark}

We subjected \textsc{ADP++} and the DS-Score to an exhaustive
benchmark spanning 21 complex topological manifolds across sample
sizes $N \in [\num{800}, \num{10000}]$ and ambient dimensions $D \in
[\num{2}, \num{100}]$:
\begin{itemize}[noitemsep,topsep=2pt]
  \item Non-convex manifolds: Concentric Rings, Two Moons, Swiss
    Roll, S-Curve, Torus ($D = \num{3}$), 4D/8D/16D Hyperspheres, Trefoil
    Knot, Mobius Strip.
  \item Extreme densities and dimensionality: Cauchy heavy-tailed
    distributions, Extreme Aspect Ratio Ellipsoids ($\num{100}:\num{1}$), and
    High-Dimensional Gaussian ($D = \num{100}$).
\end{itemize}

The draft proposes checking the following properties; its numerical
outcomes have not been verified:
\begin{enumerate}[noitemsep]
  \item \textbf{Non-Negativity:} $S_{\mathrm{PAk}} \ge 0$
    identically across all realizations and dimensionalities.
  \item \textbf{Statistical Scale Invariance:} As sample size
    increased from $N = \num{800}$ to $N = \num{10000}$, pairwise $Z$-scores
    need not scale as $\sqrt N$. Capped ranks and changing geometry
    prevent this conclusion from the central limit theorem alone;
    growth must be evaluated under a specified experiment.
  \item \textbf{Correct Topography Discovery:} \textsc{ADP++}
    correctly identified $K = \num{1}$ mode for all single-manifold
    topologies, while isolating true multi-component topologies with
    Adjusted Rand Index $\mathrm{ARI} \ge \num{0.98}$. Single-core
    execution times ranged between $\qty{0.2}{\second}$ and
    $\qty{1.6}{\second}$ per $\num{1e4}$ points.
\end{enumerate}

\subsection{Seismic Swarm Benchmark: Synthetic Scenarios}
\label{subsec:seismic_benchmark_results}

To evaluate the complete two-stage Unified Model on seismotectonic
sequences, we generated \num{100} independent Monte Carlo realizations
across three benchmark scenarios ($T_{\mathrm{obs}} =
  \qty{365}{\day}$, $m_0 = \num{1.5}$, $\beta = \num{2.30}$, $N
\approx \num{5000}$ events):

\paragraph{Scenario A (Well-Separated Sequences).}
Comprises stationary background ($\mu_0 = \qty{2.0}{\per\day}$), an
aftershock sequence at $t = \qty{50}{\day}$ ($M_{\mathrm{main}} =
\num{5.2}$, $K = \num{0.8}$, $c = \num{0.05}$, $p = \num{1.15}$),
and an independent fluid
swarm at $t = \qty{200}{\day}$ ($A = \num{1200}$, $t_0 =
  \qty{200}{\day}$, $\nu =
\num{2.5}$, $\sigma = \num{0.8}$).
\begin{itemize}[noitemsep]
  \item \textbf{Stage 1 ($\pak$):} Correctly isolated both
    transients with zero false positives: mean $Z_{12} = \num{4.1}
    \pm \num{0.3}$.
  \item \textbf{Stage 2 (Unified EM):} Accurately recovered the
    branching ratio for the aftershock sequence ($n_{\mathrm{b}} =
    \num{0.78} \pm \num{0.04}$) and identified the swarm as
    fluid-driven ($n_{\mathrm{b}} = \num{0.18} \pm \num{0.05}$,
    envelope fraction $\bar{\eta} = \num{0.82}$).
  \item Overall clustering accuracy: $\mathrm{ARI} = \num{0.97} \pm
    \num{0.02}$.
\end{itemize}

Under the unnormalized kernel in \cref{eq:clustering_omori_kernel},
these parameters give $Kc^{1-p}/(p-1)\approx8.36$, hence
$n_{\mathrm b}>8.36$ for $0<\alpha<\beta$. This is incompatible
with the draft value $0.78$. Kernel normalization and parameters
must be reconciled before using this scenario as a benchmark.

\paragraph{Scenario B (Overlapping Aftershock and Swarm).}
A fluid swarm initiates at $t = \qty{80}{\day}$ directly within the
active aftershock zone of the $t = \qty{50}{\day}$ mainshock.
\begin{itemize}[noitemsep]
  \item Standard windowing and pure \ac{ETAS} collapsed both
    sequences into a single mainshock cascade, artificially
    inflating the \ac{ETAS} branching ratio to an explosive
    $n_{\mathrm{b}} = \num{1.12} \pm \num{0.08}$.
  \item The Unified Model resolved the saddle barrier ($Z =
    \num{1.6} \pm \num{0.4}$) and correctly decoupled the three-way
    responsibilities: $\mathrm{ARI} = \num{0.82} \pm \num{0.06}$,
    representing an \textbf{\qty{+11}{\percent}} gain over baseline
    declustering
    algorithms.
\end{itemize}

The draft saddle $Z=1.6$ is below the proposed temporal threshold 3;
the stated separation needs a different explicit decision rule
and a supporting run record.

\paragraph{Scenario C (Pure Poisson Null Model).}
A stationary homogeneous Poisson process with constant rate $\mu_0 =
\qty{15.0}{\per\day}$ ($N \approx \num{5475}$ events).
\begin{itemize}[noitemsep]
  \item In \qty{98}{\percent} of realizations, the joint objective
    $\mathcal{J}(K)$ was reported to select $K^*=1$, but $K$ counts
    envelopes, so the no-envelope null requires $K=0$. This draft
    outcome does not verify false-discovery control, and ARI does
    not establish calibrated type-I error.
\end{itemize}

% ==============================================================================
% SECTION 9: DISCUSSION, LIMITATIONS, AND FUTURE DIRECTIONS
% ==============================================================================
\section{Discussion, Limitations, and Future Directions}
\label{sec:clustering_discussion}

\subsection{Summary of Accomplishments}
\label{subsec:discussion_accomplishments}

The chapter distinguishes implemented clustering operations from
ideal sampling results and proposed temporal models:
\begin{itemize}[noitemsep,topsep=2pt]
  \item Synthesised stochastic point processes, grand canonical
    thermodynamics, and information geometry, proving that
    log-density coordinates have constant Fisher metric for fixed-rank
    stopping volumes, not fixed-volume Poisson sampling.
  \item Derived fixed-rank density identities and independent-ratio
    TWO-NN limits, distinguishing them from adaptive selection and
    operational uncertainty formulas.
  \item Proposed a Unified Generative Model and joint MDL--PAk objective;
    physical attribution, identifiability, and consistency remain open.
  \item Engineered the vectorized \textsc{ADP++} clustering engine,
    proving pointer-jumping convergence for a fixed ascent forest.
    General scalar parity and draft speedups remain unverified.
  \item Reconciled PAk-DS scoring with source, related its contrasts
    to ideal persistence, and bounded the statistical and Kramers analogies.
\end{itemize}

\subsection{Limitations and Operational Boundary Conditions}
\label{subsec:discussion_limitations}

Despite these mathematical advantages, three boundary conditions
govern practical deployment:
\begin{enumerate}[noitemsep,topsep=2pt]
  \item \textbf{Short-Term Incompleteness After Major Ruptures:}
    Immediately following large earthquakes ($m \ge 6.0$), severe
    waveform overlap leads to temporary instrumentation blindness.
    Although the $c$ parameter in the Omori kernel regularizes this
    dead time, the non-parametric \pak{} estimator requires local completeness.
  \item \textbf{Boundary Truncation Effects:} Hypocentres located
    near the geographic margins of regional seismic networks suffer
    boundary volume distortion. Correcting for network convex-hull
    truncation requires boundary-corrected ball measures
    $\omega_d(\mathbf{x}_i)$.
  \item \textbf{Spatial Anisotropy:} The current \pak{} metric ball
    assumes isotropic Riemannian metrics locally. In fault zones
    with extreme fault-normal versus fault-parallel permeability
    contrasts ($k_\parallel / k_\perp \sim 10^3$), incorporating
    anisotropic metric tensors $g_{ij}(\mathbf{x})$ represents a
    natural extension.
\end{enumerate}

\subsection{Future Research Directions}
\label{subsec:discussion_future}

Three major research avenues extend naturally from this work:
\begin{itemize}[noitemsep,topsep=2pt]
  \item \textbf{Full 4D Spatiotemporal Manifold Clustering:}
    Extending the \textsc{ADP++} graph from temporal projections to
    unified $(x, y, z, t)$ spacetime, incorporating travel-time
    metric tensors derived from regional 3D velocity models.
  \item \textbf{Streaming Online Clustering:} Formulating an
    incremental, sliding-window \textsc{ADP++} engine for real-time
    seismic monitoring, alerting civil protection authorities to
    accelerating fluid-driven swarm migration in near-real-time.
  \item \textbf{Integration with Deep Phase Pickers:} Coupling the
    probabilistic cluster responsibilities $w_i^{(\mathrm{bg})},
    w_i^{(\mathrm{env})}, w_i^{(\mathrm{trig})}$ directly into the
    loss functions of fine-tuned phase pickers
    (\cref{chap:computational_pipeline,chap:training}), establishing
    an end-to-end data-centric
    pipeline from continuous waveforms to physical sequence interpretation.
\end{itemize}
